%% SISPAD 2026 — figures + captions, kept separate from the main abstract
%% file so captions can be iterated independently and dropped into
%% papers/sispad2026_abstract.tex via %% SISPAD 2026 — figures + captions, kept separate from the main abstract
%% file so captions can be iterated independently and dropped into
%% papers/sispad2026_abstract.tex via %% SISPAD 2026 — figures + captions, kept separate from the main abstract
%% file so captions can be iterated independently and dropped into
%% papers/sispad2026_abstract.tex via %% SISPAD 2026 — figures + captions, kept separate from the main abstract
%% file so captions can be iterated independently and dropped into
%% papers/sispad2026_abstract.tex via \input{sispad2026_figures.tex}.
%%
%% Each caption is written to be fully self-contained: an expert in the
%% field who reads only the paper title and this figure + caption should
%% be able to tell what was done, how, and what was shown.
%%
%% Figure files (to be generated from results/2026-04-09/*.npz by
%% run/plot_sispad_figs.py):
%%
%%   figs/fig1_device.pdf     — Device stack, band diagram, chi(z), rank-1 loop
%%   figs/fig2_patil.pdf      — Patil 1D benchmark + conservation residual inset
%%   figs/fig3_iv.pdf         — I–V of Baseline/Mol_A/Mol_B/Mol_AB on ZnO/MgZnO
%%   figs/fig4_d2i.pdf        — Analytic d2I/dV2 for all 4 molecules
%%   figs/fig5_discrimination.pdf — Baseline-subtracted discrimination metric
%%
%% Usage in sispad2026_abstract.tex:
%%   \input{sispad2026_figures}
%% or drop the individual \begin{figure}…\end{figure} blocks inline.
%% ----------------------------------------------------------------------

% --- Fig. 1 ----------------------------------------------------------------
\begin{figure}[tbp]
\centerline{\includegraphics[width=0.95\columnwidth]{figs/fig1_device.pdf}}
\caption{\textbf{Device and methodology at a glance.}
(a)~Layer stack of the asymmetric
ZnO/Mg\textsubscript{0.3}Zn\textsubscript{0.7}O RTD: 10~nm $n^+$-ZnO emitter
contact, 3~nm Mg\textsubscript{0.3}Zn\textsubscript{0.7}O emitter barrier
(host of the sensing molecule), 3~nm ZnO quantum well, 1.5~nm
Mg\textsubscript{0.3}Zn\textsubscript{0.7}O collector barrier, and 10~nm
$n^+$-ZnO collector. Conduction-band offset $\Delta E_c = 0.47$~eV,
effective mass $m^* = 0.28\,m_0$. The transverse area is
$A = 10\,\mu\mathrm{m} \times 10\,\mu\mathrm{m}$, entering the current
expression as a Datta ``Fix-$A$'' prefactor.
(b)~Conduction-band profile at $V = 0$ and at $V = V_{\mathrm{res}}$, showing
the quantized well level $E_{\mathrm{res}}$ crossing $E_F$ from above to
give negative differential resistance.
(c)~Molecular form factor $\chi(z)$, a Gaussian of width
$\sigma_{\mathrm{mol}} = 3$~\AA\ centered at the emitter-barrier midpoint,
defining the rank-1 electron--vibron coupling
$M = D_0\,\chi(z)\,|u\rangle\langle u|$.
(d)~Schematic of the rank-1 projected SCBA loop: a single
scalar-in-mode-space, 1D-in-$z$ Dyson equation dressed self-consistently by
the phonon self-energies $\tilde\sigma^{R,<,>}$. This replaces the full
nine-mode coupled SCBA with machine-precision accuracy whenever
$\sigma_{\mathrm{mol}} k_\perp \ll 1$.}
\label{fig:device}
\end{figure}

% --- Fig. 2 ----------------------------------------------------------------
\begin{figure}[tbp]
\centerline{\includegraphics[width=0.95\columnwidth]{figs/fig2_patil.pdf}}
\caption{\textbf{Validation against the Patil 1D benchmark (Tier-1 $\gamma$,
$\delta$).}
Right-contact current $I_R(V)$ of Patil's 1D symmetric RTD computed with
the rank-1 Keldysh wrapper (red circles) overlaid on the Patil 1D reference
from the published MATLAB code \texttt{rtd2modes\_1d.m} (black line). The
rank-1 solver reproduces the NDR peak to 0.55\%:
$I_R = 5.667 \times 10^{-10}$~A (rank-1) versus
$5.698 \times 10^{-10}$~A (reference) at $V = 0.336$~V. The orange
triangles show $-I_L$ on the same axes; the two traces would be exactly
coincident for a perfectly converged SCBA solution.
Inset: the per-bias current-conservation residual
$|I_L + I_R| / \max(|I_L|,|I_R|)$ in log scale. Exposing both contacts lets
us read $\delta$ directly from the Meir--Wingreen trace without any
auxiliary bookkeeping. All parameters (device, mesh, bias convention,
phonon modes, coupling, temperature) are identical to the reference.}
\label{fig:patil}
\end{figure}

% --- Fig. 3 ----------------------------------------------------------------
\begin{figure}[tbp]
\centerline{\includegraphics[width=0.95\columnwidth]{figs/fig3_iv.pdf}}
\caption{\textbf{Current--voltage characteristics of the ZnO/MgZnO primary
stack at 300~K.}
Rank-1 projected SCBA sweep (\texttt{max\_iter}$=10$, mix$=0.4$) on the
asymmetric ZnO/Mg\textsubscript{0.3}Zn\textsubscript{0.7}O RTD for the
four test species: Baseline (blank reference, bulk ZnO LO phonon at 72~meV
only, no molecule), Mol\_A (single vibrational mode $\hbar\omega_A =
100$~meV), Mol\_B ($\hbar\omega_B = 180$~meV), and Mol\_AB (both modes).
All molecules couple with $D_0^2 = 0.1\,\mathrm{eV}^2$ through a Gaussian
form factor of width 3~\AA\ centered on the emitter-barrier midpoint. Each
trace shows a clean NDR feature at $V \approx V_{\mathrm{res}}$; the phonon
sidebands pull current into the sub-resonance region, shifting the turn-on
and broadening the NDR step in a mode-specific way. Bias grid: 201~points
over 0--400~mV. Energy grid: $dE = 2$~meV over $[-0.25, +0.5]$~eV.
$E_F = 20$~meV. No numerical differentiation is applied to $I(V)$
anywhere in this work; the $d^2I/dV^2$ in Fig.~\ref{fig:d2i} comes from
the closed-form formula~(\ref{eq:d2i}) evaluated on the same converged
Keldysh state used to build this figure.}
\label{fig:iv}
\end{figure}

% --- Fig. 4 ----------------------------------------------------------------
\begin{figure}[tbp]
\centerline{\includegraphics[width=0.95\columnwidth]{figs/fig4_d2i.pdf}}
\caption{\textbf{Analytic $d^2I_R/dV^2$ --- vibrational fingerprints at
300~K.}
The second derivative is computed in closed form directly from the
converged Keldysh state at each bias via
$d^2I_R/dV^2 = (e^2/h)\cdot(1/4)\int dE\,[f_L''(E,V)\,T_{RL}(E) +
f_R''(E,V)\,(T_{RR}(E) - T_{RA}(E))]$, with the kernels
$T_{RL}, T_{RR}, T_{RA}$ built once from
$(G^R,\Gamma_L,\Gamma_R)$ at each bias --- no numerical $V$-grid
differentiation. At $D_0=0$ the formula reduces to the Landauer coherent
expression by the optical theorem
$T_{RA} = T_{RL} + T_{RR}$ (unit-tested to $10^{-6}$).
Mol\_A shows a peak at $V \approx V_{\mathrm{res}} + \hbar\omega_A/e \approx
187$~mV; Mol\_B at $V \approx V_{\mathrm{res}} + \hbar\omega_B/e \approx
267$~mV; Mol\_AB shows both features simultaneously. Baseline has a shared
feature at $\sim 162$~mV from the 72~meV bulk LO phonon. Red ticks mark the
predicted sideband biases: these are \emph{not fits}.}
\label{fig:d2i}
\end{figure}

% --- Fig. 5 ----------------------------------------------------------------
\begin{figure}[tbp]
\centerline{\includegraphics[width=0.95\columnwidth]{figs/fig5_discrimination.pdf}}
\caption{\textbf{Discrimination metric
$\Delta D(V) = d^2I_{\mathrm{mol}}/dV^2 - d^2I_{\mathrm{Baseline}}/dV^2$.}
Baseline-subtracted second-derivative spectra isolate the molecular IETS
contribution from the bulk-phonon and coherent backgrounds. The 100~meV
and 180~meV modes produce non-overlapping features separated by 80~mV,
well outside the thermal broadening $3k_BT/e \approx 78$~mV at 300~K. For
Mol\_AB the two features add linearly ($\Delta D_{AB}(V) \approx
\Delta D_A(V) + \Delta D_B(V)$ within numerical noise), a direct
demonstration of the closed-form analytic second-derivative (\ref{eq:d2i})
acting on a mixed analyte. The peak-to-peak values
$|\Delta D|_{\max}$ for each single-mode species are the figures of merit
for pixel-level discrimination. No lock-in or smoothing of any kind is
applied: the signal here is the direct output of the closed-form
(\ref{eq:d2i}) minus its Baseline evaluation.}
\label{fig:disc}
\end{figure}
.
%%
%% Each caption is written to be fully self-contained: an expert in the
%% field who reads only the paper title and this figure + caption should
%% be able to tell what was done, how, and what was shown.
%%
%% Figure files (to be generated from results/2026-04-09/*.npz by
%% run/plot_sispad_figs.py):
%%
%%   figs/fig1_device.pdf     — Device stack, band diagram, chi(z), rank-1 loop
%%   figs/fig2_patil.pdf      — Patil 1D benchmark + conservation residual inset
%%   figs/fig3_iv.pdf         — I–V of Baseline/Mol_A/Mol_B/Mol_AB on ZnO/MgZnO
%%   figs/fig4_d2i.pdf        — Analytic d2I/dV2 for all 4 molecules
%%   figs/fig5_discrimination.pdf — Baseline-subtracted discrimination metric
%%
%% Usage in sispad2026_abstract.tex:
%%   %% SISPAD 2026 — figures + captions, kept separate from the main abstract
%% file so captions can be iterated independently and dropped into
%% papers/sispad2026_abstract.tex via \input{sispad2026_figures.tex}.
%%
%% Each caption is written to be fully self-contained: an expert in the
%% field who reads only the paper title and this figure + caption should
%% be able to tell what was done, how, and what was shown.
%%
%% Figure files (to be generated from results/2026-04-09/*.npz by
%% run/plot_sispad_figs.py):
%%
%%   figs/fig1_device.pdf     — Device stack, band diagram, chi(z), rank-1 loop
%%   figs/fig2_patil.pdf      — Patil 1D benchmark + conservation residual inset
%%   figs/fig3_iv.pdf         — I–V of Baseline/Mol_A/Mol_B/Mol_AB on ZnO/MgZnO
%%   figs/fig4_d2i.pdf        — Analytic d2I/dV2 for all 4 molecules
%%   figs/fig5_discrimination.pdf — Baseline-subtracted discrimination metric
%%
%% Usage in sispad2026_abstract.tex:
%%   \input{sispad2026_figures}
%% or drop the individual \begin{figure}…\end{figure} blocks inline.
%% ----------------------------------------------------------------------

% --- Fig. 1 ----------------------------------------------------------------
\begin{figure}[tbp]
\centerline{\includegraphics[width=0.95\columnwidth]{figs/fig1_device.pdf}}
\caption{\textbf{Device and methodology at a glance.}
(a)~Layer stack of the asymmetric
ZnO/Mg\textsubscript{0.3}Zn\textsubscript{0.7}O RTD: 10~nm $n^+$-ZnO emitter
contact, 3~nm Mg\textsubscript{0.3}Zn\textsubscript{0.7}O emitter barrier
(host of the sensing molecule), 3~nm ZnO quantum well, 1.5~nm
Mg\textsubscript{0.3}Zn\textsubscript{0.7}O collector barrier, and 10~nm
$n^+$-ZnO collector. Conduction-band offset $\Delta E_c = 0.47$~eV,
effective mass $m^* = 0.28\,m_0$. The transverse area is
$A = 10\,\mu\mathrm{m} \times 10\,\mu\mathrm{m}$, entering the current
expression as a Datta ``Fix-$A$'' prefactor.
(b)~Conduction-band profile at $V = 0$ and at $V = V_{\mathrm{res}}$, showing
the quantized well level $E_{\mathrm{res}}$ crossing $E_F$ from above to
give negative differential resistance.
(c)~Molecular form factor $\chi(z)$, a Gaussian of width
$\sigma_{\mathrm{mol}} = 3$~\AA\ centered at the emitter-barrier midpoint,
defining the rank-1 electron--vibron coupling
$M = D_0\,\chi(z)\,|u\rangle\langle u|$.
(d)~Schematic of the rank-1 projected SCBA loop: a single
scalar-in-mode-space, 1D-in-$z$ Dyson equation dressed self-consistently by
the phonon self-energies $\tilde\sigma^{R,<,>}$. This replaces the full
nine-mode coupled SCBA with machine-precision accuracy whenever
$\sigma_{\mathrm{mol}} k_\perp \ll 1$.}
\label{fig:device}
\end{figure}

% --- Fig. 2 ----------------------------------------------------------------
\begin{figure}[tbp]
\centerline{\includegraphics[width=0.95\columnwidth]{figs/fig2_patil.pdf}}
\caption{\textbf{Validation against the Patil 1D benchmark (Tier-1 $\gamma$,
$\delta$).}
Right-contact current $I_R(V)$ of Patil's 1D symmetric RTD computed with
the rank-1 Keldysh wrapper (red circles) overlaid on the Patil 1D reference
from the published MATLAB code \texttt{rtd2modes\_1d.m} (black line). The
rank-1 solver reproduces the NDR peak to 0.55\%:
$I_R = 5.667 \times 10^{-10}$~A (rank-1) versus
$5.698 \times 10^{-10}$~A (reference) at $V = 0.336$~V. The orange
triangles show $-I_L$ on the same axes; the two traces would be exactly
coincident for a perfectly converged SCBA solution.
Inset: the per-bias current-conservation residual
$|I_L + I_R| / \max(|I_L|,|I_R|)$ in log scale. Exposing both contacts lets
us read $\delta$ directly from the Meir--Wingreen trace without any
auxiliary bookkeeping. All parameters (device, mesh, bias convention,
phonon modes, coupling, temperature) are identical to the reference.}
\label{fig:patil}
\end{figure}

% --- Fig. 3 ----------------------------------------------------------------
\begin{figure}[tbp]
\centerline{\includegraphics[width=0.95\columnwidth]{figs/fig3_iv.pdf}}
\caption{\textbf{Current--voltage characteristics of the ZnO/MgZnO primary
stack at 300~K.}
Rank-1 projected SCBA sweep (\texttt{max\_iter}$=10$, mix$=0.4$) on the
asymmetric ZnO/Mg\textsubscript{0.3}Zn\textsubscript{0.7}O RTD for the
four test species: Baseline (blank reference, bulk ZnO LO phonon at 72~meV
only, no molecule), Mol\_A (single vibrational mode $\hbar\omega_A =
100$~meV), Mol\_B ($\hbar\omega_B = 180$~meV), and Mol\_AB (both modes).
All molecules couple with $D_0^2 = 0.1\,\mathrm{eV}^2$ through a Gaussian
form factor of width 3~\AA\ centered on the emitter-barrier midpoint. Each
trace shows a clean NDR feature at $V \approx V_{\mathrm{res}}$; the phonon
sidebands pull current into the sub-resonance region, shifting the turn-on
and broadening the NDR step in a mode-specific way. Bias grid: 201~points
over 0--400~mV. Energy grid: $dE = 2$~meV over $[-0.25, +0.5]$~eV.
$E_F = 20$~meV. No numerical differentiation is applied to $I(V)$
anywhere in this work; the $d^2I/dV^2$ in Fig.~\ref{fig:d2i} comes from
the closed-form formula~(\ref{eq:d2i}) evaluated on the same converged
Keldysh state used to build this figure.}
\label{fig:iv}
\end{figure}

% --- Fig. 4 ----------------------------------------------------------------
\begin{figure}[tbp]
\centerline{\includegraphics[width=0.95\columnwidth]{figs/fig4_d2i.pdf}}
\caption{\textbf{Analytic $d^2I_R/dV^2$ --- vibrational fingerprints at
300~K.}
The second derivative is computed in closed form directly from the
converged Keldysh state at each bias via
$d^2I_R/dV^2 = (e^2/h)\cdot(1/4)\int dE\,[f_L''(E,V)\,T_{RL}(E) +
f_R''(E,V)\,(T_{RR}(E) - T_{RA}(E))]$, with the kernels
$T_{RL}, T_{RR}, T_{RA}$ built once from
$(G^R,\Gamma_L,\Gamma_R)$ at each bias --- no numerical $V$-grid
differentiation. At $D_0=0$ the formula reduces to the Landauer coherent
expression by the optical theorem
$T_{RA} = T_{RL} + T_{RR}$ (unit-tested to $10^{-6}$).
Mol\_A shows a peak at $V \approx V_{\mathrm{res}} + \hbar\omega_A/e \approx
187$~mV; Mol\_B at $V \approx V_{\mathrm{res}} + \hbar\omega_B/e \approx
267$~mV; Mol\_AB shows both features simultaneously. Baseline has a shared
feature at $\sim 162$~mV from the 72~meV bulk LO phonon. Red ticks mark the
predicted sideband biases: these are \emph{not fits}.}
\label{fig:d2i}
\end{figure}

% --- Fig. 5 ----------------------------------------------------------------
\begin{figure}[tbp]
\centerline{\includegraphics[width=0.95\columnwidth]{figs/fig5_discrimination.pdf}}
\caption{\textbf{Discrimination metric
$\Delta D(V) = d^2I_{\mathrm{mol}}/dV^2 - d^2I_{\mathrm{Baseline}}/dV^2$.}
Baseline-subtracted second-derivative spectra isolate the molecular IETS
contribution from the bulk-phonon and coherent backgrounds. The 100~meV
and 180~meV modes produce non-overlapping features separated by 80~mV,
well outside the thermal broadening $3k_BT/e \approx 78$~mV at 300~K. For
Mol\_AB the two features add linearly ($\Delta D_{AB}(V) \approx
\Delta D_A(V) + \Delta D_B(V)$ within numerical noise), a direct
demonstration of the closed-form analytic second-derivative (\ref{eq:d2i})
acting on a mixed analyte. The peak-to-peak values
$|\Delta D|_{\max}$ for each single-mode species are the figures of merit
for pixel-level discrimination. No lock-in or smoothing of any kind is
applied: the signal here is the direct output of the closed-form
(\ref{eq:d2i}) minus its Baseline evaluation.}
\label{fig:disc}
\end{figure}

%% or drop the individual \begin{figure}…\end{figure} blocks inline.
%% ----------------------------------------------------------------------

% --- Fig. 1 ----------------------------------------------------------------
\begin{figure}[tbp]
\centerline{\includegraphics[width=0.95\columnwidth]{figs/fig1_device.pdf}}
\caption{\textbf{Device and methodology at a glance.}
(a)~Layer stack of the asymmetric
ZnO/Mg\textsubscript{0.3}Zn\textsubscript{0.7}O RTD: 10~nm $n^+$-ZnO emitter
contact, 3~nm Mg\textsubscript{0.3}Zn\textsubscript{0.7}O emitter barrier
(host of the sensing molecule), 3~nm ZnO quantum well, 1.5~nm
Mg\textsubscript{0.3}Zn\textsubscript{0.7}O collector barrier, and 10~nm
$n^+$-ZnO collector. Conduction-band offset $\Delta E_c = 0.47$~eV,
effective mass $m^* = 0.28\,m_0$. The transverse area is
$A = 10\,\mu\mathrm{m} \times 10\,\mu\mathrm{m}$, entering the current
expression as a Datta ``Fix-$A$'' prefactor.
(b)~Conduction-band profile at $V = 0$ and at $V = V_{\mathrm{res}}$, showing
the quantized well level $E_{\mathrm{res}}$ crossing $E_F$ from above to
give negative differential resistance.
(c)~Molecular form factor $\chi(z)$, a Gaussian of width
$\sigma_{\mathrm{mol}} = 3$~\AA\ centered at the emitter-barrier midpoint,
defining the rank-1 electron--vibron coupling
$M = D_0\,\chi(z)\,|u\rangle\langle u|$.
(d)~Schematic of the rank-1 projected SCBA loop: a single
scalar-in-mode-space, 1D-in-$z$ Dyson equation dressed self-consistently by
the phonon self-energies $\tilde\sigma^{R,<,>}$. This replaces the full
nine-mode coupled SCBA with machine-precision accuracy whenever
$\sigma_{\mathrm{mol}} k_\perp \ll 1$.}
\label{fig:device}
\end{figure}

% --- Fig. 2 ----------------------------------------------------------------
\begin{figure}[tbp]
\centerline{\includegraphics[width=0.95\columnwidth]{figs/fig2_patil.pdf}}
\caption{\textbf{Validation against the Patil 1D benchmark (Tier-1 $\gamma$,
$\delta$).}
Right-contact current $I_R(V)$ of Patil's 1D symmetric RTD computed with
the rank-1 Keldysh wrapper (red circles) overlaid on the Patil 1D reference
from the published MATLAB code \texttt{rtd2modes\_1d.m} (black line). The
rank-1 solver reproduces the NDR peak to 0.55\%:
$I_R = 5.667 \times 10^{-10}$~A (rank-1) versus
$5.698 \times 10^{-10}$~A (reference) at $V = 0.336$~V. The orange
triangles show $-I_L$ on the same axes; the two traces would be exactly
coincident for a perfectly converged SCBA solution.
Inset: the per-bias current-conservation residual
$|I_L + I_R| / \max(|I_L|,|I_R|)$ in log scale. Exposing both contacts lets
us read $\delta$ directly from the Meir--Wingreen trace without any
auxiliary bookkeeping. All parameters (device, mesh, bias convention,
phonon modes, coupling, temperature) are identical to the reference.}
\label{fig:patil}
\end{figure}

% --- Fig. 3 ----------------------------------------------------------------
\begin{figure}[tbp]
\centerline{\includegraphics[width=0.95\columnwidth]{figs/fig3_iv.pdf}}
\caption{\textbf{Current--voltage characteristics of the ZnO/MgZnO primary
stack at 300~K.}
Rank-1 projected SCBA sweep (\texttt{max\_iter}$=10$, mix$=0.4$) on the
asymmetric ZnO/Mg\textsubscript{0.3}Zn\textsubscript{0.7}O RTD for the
four test species: Baseline (blank reference, bulk ZnO LO phonon at 72~meV
only, no molecule), Mol\_A (single vibrational mode $\hbar\omega_A =
100$~meV), Mol\_B ($\hbar\omega_B = 180$~meV), and Mol\_AB (both modes).
All molecules couple with $D_0^2 = 0.1\,\mathrm{eV}^2$ through a Gaussian
form factor of width 3~\AA\ centered on the emitter-barrier midpoint. Each
trace shows a clean NDR feature at $V \approx V_{\mathrm{res}}$; the phonon
sidebands pull current into the sub-resonance region, shifting the turn-on
and broadening the NDR step in a mode-specific way. Bias grid: 201~points
over 0--400~mV. Energy grid: $dE = 2$~meV over $[-0.25, +0.5]$~eV.
$E_F = 20$~meV. No numerical differentiation is applied to $I(V)$
anywhere in this work; the $d^2I/dV^2$ in Fig.~\ref{fig:d2i} comes from
the closed-form formula~(\ref{eq:d2i}) evaluated on the same converged
Keldysh state used to build this figure.}
\label{fig:iv}
\end{figure}

% --- Fig. 4 ----------------------------------------------------------------
\begin{figure}[tbp]
\centerline{\includegraphics[width=0.95\columnwidth]{figs/fig4_d2i.pdf}}
\caption{\textbf{Analytic $d^2I_R/dV^2$ --- vibrational fingerprints at
300~K.}
The second derivative is computed in closed form directly from the
converged Keldysh state at each bias via
$d^2I_R/dV^2 = (e^2/h)\cdot(1/4)\int dE\,[f_L''(E,V)\,T_{RL}(E) +
f_R''(E,V)\,(T_{RR}(E) - T_{RA}(E))]$, with the kernels
$T_{RL}, T_{RR}, T_{RA}$ built once from
$(G^R,\Gamma_L,\Gamma_R)$ at each bias --- no numerical $V$-grid
differentiation. At $D_0=0$ the formula reduces to the Landauer coherent
expression by the optical theorem
$T_{RA} = T_{RL} + T_{RR}$ (unit-tested to $10^{-6}$).
Mol\_A shows a peak at $V \approx V_{\mathrm{res}} + \hbar\omega_A/e \approx
187$~mV; Mol\_B at $V \approx V_{\mathrm{res}} + \hbar\omega_B/e \approx
267$~mV; Mol\_AB shows both features simultaneously. Baseline has a shared
feature at $\sim 162$~mV from the 72~meV bulk LO phonon. Red ticks mark the
predicted sideband biases: these are \emph{not fits}.}
\label{fig:d2i}
\end{figure}

% --- Fig. 5 ----------------------------------------------------------------
\begin{figure}[tbp]
\centerline{\includegraphics[width=0.95\columnwidth]{figs/fig5_discrimination.pdf}}
\caption{\textbf{Discrimination metric
$\Delta D(V) = d^2I_{\mathrm{mol}}/dV^2 - d^2I_{\mathrm{Baseline}}/dV^2$.}
Baseline-subtracted second-derivative spectra isolate the molecular IETS
contribution from the bulk-phonon and coherent backgrounds. The 100~meV
and 180~meV modes produce non-overlapping features separated by 80~mV,
well outside the thermal broadening $3k_BT/e \approx 78$~mV at 300~K. For
Mol\_AB the two features add linearly ($\Delta D_{AB}(V) \approx
\Delta D_A(V) + \Delta D_B(V)$ within numerical noise), a direct
demonstration of the closed-form analytic second-derivative (\ref{eq:d2i})
acting on a mixed analyte. The peak-to-peak values
$|\Delta D|_{\max}$ for each single-mode species are the figures of merit
for pixel-level discrimination. No lock-in or smoothing of any kind is
applied: the signal here is the direct output of the closed-form
(\ref{eq:d2i}) minus its Baseline evaluation.}
\label{fig:disc}
\end{figure}
.
%%
%% Each caption is written to be fully self-contained: an expert in the
%% field who reads only the paper title and this figure + caption should
%% be able to tell what was done, how, and what was shown.
%%
%% Figure files (to be generated from results/2026-04-09/*.npz by
%% run/plot_sispad_figs.py):
%%
%%   figs/fig1_device.pdf     — Device stack, band diagram, chi(z), rank-1 loop
%%   figs/fig2_patil.pdf      — Patil 1D benchmark + conservation residual inset
%%   figs/fig3_iv.pdf         — I–V of Baseline/Mol_A/Mol_B/Mol_AB on ZnO/MgZnO
%%   figs/fig4_d2i.pdf        — Analytic d2I/dV2 for all 4 molecules
%%   figs/fig5_discrimination.pdf — Baseline-subtracted discrimination metric
%%
%% Usage in sispad2026_abstract.tex:
%%   %% SISPAD 2026 — figures + captions, kept separate from the main abstract
%% file so captions can be iterated independently and dropped into
%% papers/sispad2026_abstract.tex via %% SISPAD 2026 — figures + captions, kept separate from the main abstract
%% file so captions can be iterated independently and dropped into
%% papers/sispad2026_abstract.tex via \input{sispad2026_figures.tex}.
%%
%% Each caption is written to be fully self-contained: an expert in the
%% field who reads only the paper title and this figure + caption should
%% be able to tell what was done, how, and what was shown.
%%
%% Figure files (to be generated from results/2026-04-09/*.npz by
%% run/plot_sispad_figs.py):
%%
%%   figs/fig1_device.pdf     — Device stack, band diagram, chi(z), rank-1 loop
%%   figs/fig2_patil.pdf      — Patil 1D benchmark + conservation residual inset
%%   figs/fig3_iv.pdf         — I–V of Baseline/Mol_A/Mol_B/Mol_AB on ZnO/MgZnO
%%   figs/fig4_d2i.pdf        — Analytic d2I/dV2 for all 4 molecules
%%   figs/fig5_discrimination.pdf — Baseline-subtracted discrimination metric
%%
%% Usage in sispad2026_abstract.tex:
%%   \input{sispad2026_figures}
%% or drop the individual \begin{figure}…\end{figure} blocks inline.
%% ----------------------------------------------------------------------

% --- Fig. 1 ----------------------------------------------------------------
\begin{figure}[tbp]
\centerline{\includegraphics[width=0.95\columnwidth]{figs/fig1_device.pdf}}
\caption{\textbf{Device and methodology at a glance.}
(a)~Layer stack of the asymmetric
ZnO/Mg\textsubscript{0.3}Zn\textsubscript{0.7}O RTD: 10~nm $n^+$-ZnO emitter
contact, 3~nm Mg\textsubscript{0.3}Zn\textsubscript{0.7}O emitter barrier
(host of the sensing molecule), 3~nm ZnO quantum well, 1.5~nm
Mg\textsubscript{0.3}Zn\textsubscript{0.7}O collector barrier, and 10~nm
$n^+$-ZnO collector. Conduction-band offset $\Delta E_c = 0.47$~eV,
effective mass $m^* = 0.28\,m_0$. The transverse area is
$A = 10\,\mu\mathrm{m} \times 10\,\mu\mathrm{m}$, entering the current
expression as a Datta ``Fix-$A$'' prefactor.
(b)~Conduction-band profile at $V = 0$ and at $V = V_{\mathrm{res}}$, showing
the quantized well level $E_{\mathrm{res}}$ crossing $E_F$ from above to
give negative differential resistance.
(c)~Molecular form factor $\chi(z)$, a Gaussian of width
$\sigma_{\mathrm{mol}} = 3$~\AA\ centered at the emitter-barrier midpoint,
defining the rank-1 electron--vibron coupling
$M = D_0\,\chi(z)\,|u\rangle\langle u|$.
(d)~Schematic of the rank-1 projected SCBA loop: a single
scalar-in-mode-space, 1D-in-$z$ Dyson equation dressed self-consistently by
the phonon self-energies $\tilde\sigma^{R,<,>}$. This replaces the full
nine-mode coupled SCBA with machine-precision accuracy whenever
$\sigma_{\mathrm{mol}} k_\perp \ll 1$.}
\label{fig:device}
\end{figure}

% --- Fig. 2 ----------------------------------------------------------------
\begin{figure}[tbp]
\centerline{\includegraphics[width=0.95\columnwidth]{figs/fig2_patil.pdf}}
\caption{\textbf{Validation against the Patil 1D benchmark (Tier-1 $\gamma$,
$\delta$).}
Right-contact current $I_R(V)$ of Patil's 1D symmetric RTD computed with
the rank-1 Keldysh wrapper (red circles) overlaid on the Patil 1D reference
from the published MATLAB code \texttt{rtd2modes\_1d.m} (black line). The
rank-1 solver reproduces the NDR peak to 0.55\%:
$I_R = 5.667 \times 10^{-10}$~A (rank-1) versus
$5.698 \times 10^{-10}$~A (reference) at $V = 0.336$~V. The orange
triangles show $-I_L$ on the same axes; the two traces would be exactly
coincident for a perfectly converged SCBA solution.
Inset: the per-bias current-conservation residual
$|I_L + I_R| / \max(|I_L|,|I_R|)$ in log scale. Exposing both contacts lets
us read $\delta$ directly from the Meir--Wingreen trace without any
auxiliary bookkeeping. All parameters (device, mesh, bias convention,
phonon modes, coupling, temperature) are identical to the reference.}
\label{fig:patil}
\end{figure}

% --- Fig. 3 ----------------------------------------------------------------
\begin{figure}[tbp]
\centerline{\includegraphics[width=0.95\columnwidth]{figs/fig3_iv.pdf}}
\caption{\textbf{Current--voltage characteristics of the ZnO/MgZnO primary
stack at 300~K.}
Rank-1 projected SCBA sweep (\texttt{max\_iter}$=10$, mix$=0.4$) on the
asymmetric ZnO/Mg\textsubscript{0.3}Zn\textsubscript{0.7}O RTD for the
four test species: Baseline (blank reference, bulk ZnO LO phonon at 72~meV
only, no molecule), Mol\_A (single vibrational mode $\hbar\omega_A =
100$~meV), Mol\_B ($\hbar\omega_B = 180$~meV), and Mol\_AB (both modes).
All molecules couple with $D_0^2 = 0.1\,\mathrm{eV}^2$ through a Gaussian
form factor of width 3~\AA\ centered on the emitter-barrier midpoint. Each
trace shows a clean NDR feature at $V \approx V_{\mathrm{res}}$; the phonon
sidebands pull current into the sub-resonance region, shifting the turn-on
and broadening the NDR step in a mode-specific way. Bias grid: 201~points
over 0--400~mV. Energy grid: $dE = 2$~meV over $[-0.25, +0.5]$~eV.
$E_F = 20$~meV. No numerical differentiation is applied to $I(V)$
anywhere in this work; the $d^2I/dV^2$ in Fig.~\ref{fig:d2i} comes from
the closed-form formula~(\ref{eq:d2i}) evaluated on the same converged
Keldysh state used to build this figure.}
\label{fig:iv}
\end{figure}

% --- Fig. 4 ----------------------------------------------------------------
\begin{figure}[tbp]
\centerline{\includegraphics[width=0.95\columnwidth]{figs/fig4_d2i.pdf}}
\caption{\textbf{Analytic $d^2I_R/dV^2$ --- vibrational fingerprints at
300~K.}
The second derivative is computed in closed form directly from the
converged Keldysh state at each bias via
$d^2I_R/dV^2 = (e^2/h)\cdot(1/4)\int dE\,[f_L''(E,V)\,T_{RL}(E) +
f_R''(E,V)\,(T_{RR}(E) - T_{RA}(E))]$, with the kernels
$T_{RL}, T_{RR}, T_{RA}$ built once from
$(G^R,\Gamma_L,\Gamma_R)$ at each bias --- no numerical $V$-grid
differentiation. At $D_0=0$ the formula reduces to the Landauer coherent
expression by the optical theorem
$T_{RA} = T_{RL} + T_{RR}$ (unit-tested to $10^{-6}$).
Mol\_A shows a peak at $V \approx V_{\mathrm{res}} + \hbar\omega_A/e \approx
187$~mV; Mol\_B at $V \approx V_{\mathrm{res}} + \hbar\omega_B/e \approx
267$~mV; Mol\_AB shows both features simultaneously. Baseline has a shared
feature at $\sim 162$~mV from the 72~meV bulk LO phonon. Red ticks mark the
predicted sideband biases: these are \emph{not fits}.}
\label{fig:d2i}
\end{figure}

% --- Fig. 5 ----------------------------------------------------------------
\begin{figure}[tbp]
\centerline{\includegraphics[width=0.95\columnwidth]{figs/fig5_discrimination.pdf}}
\caption{\textbf{Discrimination metric
$\Delta D(V) = d^2I_{\mathrm{mol}}/dV^2 - d^2I_{\mathrm{Baseline}}/dV^2$.}
Baseline-subtracted second-derivative spectra isolate the molecular IETS
contribution from the bulk-phonon and coherent backgrounds. The 100~meV
and 180~meV modes produce non-overlapping features separated by 80~mV,
well outside the thermal broadening $3k_BT/e \approx 78$~mV at 300~K. For
Mol\_AB the two features add linearly ($\Delta D_{AB}(V) \approx
\Delta D_A(V) + \Delta D_B(V)$ within numerical noise), a direct
demonstration of the closed-form analytic second-derivative (\ref{eq:d2i})
acting on a mixed analyte. The peak-to-peak values
$|\Delta D|_{\max}$ for each single-mode species are the figures of merit
for pixel-level discrimination. No lock-in or smoothing of any kind is
applied: the signal here is the direct output of the closed-form
(\ref{eq:d2i}) minus its Baseline evaluation.}
\label{fig:disc}
\end{figure}
.
%%
%% Each caption is written to be fully self-contained: an expert in the
%% field who reads only the paper title and this figure + caption should
%% be able to tell what was done, how, and what was shown.
%%
%% Figure files (to be generated from results/2026-04-09/*.npz by
%% run/plot_sispad_figs.py):
%%
%%   figs/fig1_device.pdf     — Device stack, band diagram, chi(z), rank-1 loop
%%   figs/fig2_patil.pdf      — Patil 1D benchmark + conservation residual inset
%%   figs/fig3_iv.pdf         — I–V of Baseline/Mol_A/Mol_B/Mol_AB on ZnO/MgZnO
%%   figs/fig4_d2i.pdf        — Analytic d2I/dV2 for all 4 molecules
%%   figs/fig5_discrimination.pdf — Baseline-subtracted discrimination metric
%%
%% Usage in sispad2026_abstract.tex:
%%   %% SISPAD 2026 — figures + captions, kept separate from the main abstract
%% file so captions can be iterated independently and dropped into
%% papers/sispad2026_abstract.tex via \input{sispad2026_figures.tex}.
%%
%% Each caption is written to be fully self-contained: an expert in the
%% field who reads only the paper title and this figure + caption should
%% be able to tell what was done, how, and what was shown.
%%
%% Figure files (to be generated from results/2026-04-09/*.npz by
%% run/plot_sispad_figs.py):
%%
%%   figs/fig1_device.pdf     — Device stack, band diagram, chi(z), rank-1 loop
%%   figs/fig2_patil.pdf      — Patil 1D benchmark + conservation residual inset
%%   figs/fig3_iv.pdf         — I–V of Baseline/Mol_A/Mol_B/Mol_AB on ZnO/MgZnO
%%   figs/fig4_d2i.pdf        — Analytic d2I/dV2 for all 4 molecules
%%   figs/fig5_discrimination.pdf — Baseline-subtracted discrimination metric
%%
%% Usage in sispad2026_abstract.tex:
%%   \input{sispad2026_figures}
%% or drop the individual \begin{figure}…\end{figure} blocks inline.
%% ----------------------------------------------------------------------

% --- Fig. 1 ----------------------------------------------------------------
\begin{figure}[tbp]
\centerline{\includegraphics[width=0.95\columnwidth]{figs/fig1_device.pdf}}
\caption{\textbf{Device and methodology at a glance.}
(a)~Layer stack of the asymmetric
ZnO/Mg\textsubscript{0.3}Zn\textsubscript{0.7}O RTD: 10~nm $n^+$-ZnO emitter
contact, 3~nm Mg\textsubscript{0.3}Zn\textsubscript{0.7}O emitter barrier
(host of the sensing molecule), 3~nm ZnO quantum well, 1.5~nm
Mg\textsubscript{0.3}Zn\textsubscript{0.7}O collector barrier, and 10~nm
$n^+$-ZnO collector. Conduction-band offset $\Delta E_c = 0.47$~eV,
effective mass $m^* = 0.28\,m_0$. The transverse area is
$A = 10\,\mu\mathrm{m} \times 10\,\mu\mathrm{m}$, entering the current
expression as a Datta ``Fix-$A$'' prefactor.
(b)~Conduction-band profile at $V = 0$ and at $V = V_{\mathrm{res}}$, showing
the quantized well level $E_{\mathrm{res}}$ crossing $E_F$ from above to
give negative differential resistance.
(c)~Molecular form factor $\chi(z)$, a Gaussian of width
$\sigma_{\mathrm{mol}} = 3$~\AA\ centered at the emitter-barrier midpoint,
defining the rank-1 electron--vibron coupling
$M = D_0\,\chi(z)\,|u\rangle\langle u|$.
(d)~Schematic of the rank-1 projected SCBA loop: a single
scalar-in-mode-space, 1D-in-$z$ Dyson equation dressed self-consistently by
the phonon self-energies $\tilde\sigma^{R,<,>}$. This replaces the full
nine-mode coupled SCBA with machine-precision accuracy whenever
$\sigma_{\mathrm{mol}} k_\perp \ll 1$.}
\label{fig:device}
\end{figure}

% --- Fig. 2 ----------------------------------------------------------------
\begin{figure}[tbp]
\centerline{\includegraphics[width=0.95\columnwidth]{figs/fig2_patil.pdf}}
\caption{\textbf{Validation against the Patil 1D benchmark (Tier-1 $\gamma$,
$\delta$).}
Right-contact current $I_R(V)$ of Patil's 1D symmetric RTD computed with
the rank-1 Keldysh wrapper (red circles) overlaid on the Patil 1D reference
from the published MATLAB code \texttt{rtd2modes\_1d.m} (black line). The
rank-1 solver reproduces the NDR peak to 0.55\%:
$I_R = 5.667 \times 10^{-10}$~A (rank-1) versus
$5.698 \times 10^{-10}$~A (reference) at $V = 0.336$~V. The orange
triangles show $-I_L$ on the same axes; the two traces would be exactly
coincident for a perfectly converged SCBA solution.
Inset: the per-bias current-conservation residual
$|I_L + I_R| / \max(|I_L|,|I_R|)$ in log scale. Exposing both contacts lets
us read $\delta$ directly from the Meir--Wingreen trace without any
auxiliary bookkeeping. All parameters (device, mesh, bias convention,
phonon modes, coupling, temperature) are identical to the reference.}
\label{fig:patil}
\end{figure}

% --- Fig. 3 ----------------------------------------------------------------
\begin{figure}[tbp]
\centerline{\includegraphics[width=0.95\columnwidth]{figs/fig3_iv.pdf}}
\caption{\textbf{Current--voltage characteristics of the ZnO/MgZnO primary
stack at 300~K.}
Rank-1 projected SCBA sweep (\texttt{max\_iter}$=10$, mix$=0.4$) on the
asymmetric ZnO/Mg\textsubscript{0.3}Zn\textsubscript{0.7}O RTD for the
four test species: Baseline (blank reference, bulk ZnO LO phonon at 72~meV
only, no molecule), Mol\_A (single vibrational mode $\hbar\omega_A =
100$~meV), Mol\_B ($\hbar\omega_B = 180$~meV), and Mol\_AB (both modes).
All molecules couple with $D_0^2 = 0.1\,\mathrm{eV}^2$ through a Gaussian
form factor of width 3~\AA\ centered on the emitter-barrier midpoint. Each
trace shows a clean NDR feature at $V \approx V_{\mathrm{res}}$; the phonon
sidebands pull current into the sub-resonance region, shifting the turn-on
and broadening the NDR step in a mode-specific way. Bias grid: 201~points
over 0--400~mV. Energy grid: $dE = 2$~meV over $[-0.25, +0.5]$~eV.
$E_F = 20$~meV. No numerical differentiation is applied to $I(V)$
anywhere in this work; the $d^2I/dV^2$ in Fig.~\ref{fig:d2i} comes from
the closed-form formula~(\ref{eq:d2i}) evaluated on the same converged
Keldysh state used to build this figure.}
\label{fig:iv}
\end{figure}

% --- Fig. 4 ----------------------------------------------------------------
\begin{figure}[tbp]
\centerline{\includegraphics[width=0.95\columnwidth]{figs/fig4_d2i.pdf}}
\caption{\textbf{Analytic $d^2I_R/dV^2$ --- vibrational fingerprints at
300~K.}
The second derivative is computed in closed form directly from the
converged Keldysh state at each bias via
$d^2I_R/dV^2 = (e^2/h)\cdot(1/4)\int dE\,[f_L''(E,V)\,T_{RL}(E) +
f_R''(E,V)\,(T_{RR}(E) - T_{RA}(E))]$, with the kernels
$T_{RL}, T_{RR}, T_{RA}$ built once from
$(G^R,\Gamma_L,\Gamma_R)$ at each bias --- no numerical $V$-grid
differentiation. At $D_0=0$ the formula reduces to the Landauer coherent
expression by the optical theorem
$T_{RA} = T_{RL} + T_{RR}$ (unit-tested to $10^{-6}$).
Mol\_A shows a peak at $V \approx V_{\mathrm{res}} + \hbar\omega_A/e \approx
187$~mV; Mol\_B at $V \approx V_{\mathrm{res}} + \hbar\omega_B/e \approx
267$~mV; Mol\_AB shows both features simultaneously. Baseline has a shared
feature at $\sim 162$~mV from the 72~meV bulk LO phonon. Red ticks mark the
predicted sideband biases: these are \emph{not fits}.}
\label{fig:d2i}
\end{figure}

% --- Fig. 5 ----------------------------------------------------------------
\begin{figure}[tbp]
\centerline{\includegraphics[width=0.95\columnwidth]{figs/fig5_discrimination.pdf}}
\caption{\textbf{Discrimination metric
$\Delta D(V) = d^2I_{\mathrm{mol}}/dV^2 - d^2I_{\mathrm{Baseline}}/dV^2$.}
Baseline-subtracted second-derivative spectra isolate the molecular IETS
contribution from the bulk-phonon and coherent backgrounds. The 100~meV
and 180~meV modes produce non-overlapping features separated by 80~mV,
well outside the thermal broadening $3k_BT/e \approx 78$~mV at 300~K. For
Mol\_AB the two features add linearly ($\Delta D_{AB}(V) \approx
\Delta D_A(V) + \Delta D_B(V)$ within numerical noise), a direct
demonstration of the closed-form analytic second-derivative (\ref{eq:d2i})
acting on a mixed analyte. The peak-to-peak values
$|\Delta D|_{\max}$ for each single-mode species are the figures of merit
for pixel-level discrimination. No lock-in or smoothing of any kind is
applied: the signal here is the direct output of the closed-form
(\ref{eq:d2i}) minus its Baseline evaluation.}
\label{fig:disc}
\end{figure}

%% or drop the individual \begin{figure}…\end{figure} blocks inline.
%% ----------------------------------------------------------------------

% --- Fig. 1 ----------------------------------------------------------------
\begin{figure}[tbp]
\centerline{\includegraphics[width=0.95\columnwidth]{figs/fig1_device.pdf}}
\caption{\textbf{Device and methodology at a glance.}
(a)~Layer stack of the asymmetric
ZnO/Mg\textsubscript{0.3}Zn\textsubscript{0.7}O RTD: 10~nm $n^+$-ZnO emitter
contact, 3~nm Mg\textsubscript{0.3}Zn\textsubscript{0.7}O emitter barrier
(host of the sensing molecule), 3~nm ZnO quantum well, 1.5~nm
Mg\textsubscript{0.3}Zn\textsubscript{0.7}O collector barrier, and 10~nm
$n^+$-ZnO collector. Conduction-band offset $\Delta E_c = 0.47$~eV,
effective mass $m^* = 0.28\,m_0$. The transverse area is
$A = 10\,\mu\mathrm{m} \times 10\,\mu\mathrm{m}$, entering the current
expression as a Datta ``Fix-$A$'' prefactor.
(b)~Conduction-band profile at $V = 0$ and at $V = V_{\mathrm{res}}$, showing
the quantized well level $E_{\mathrm{res}}$ crossing $E_F$ from above to
give negative differential resistance.
(c)~Molecular form factor $\chi(z)$, a Gaussian of width
$\sigma_{\mathrm{mol}} = 3$~\AA\ centered at the emitter-barrier midpoint,
defining the rank-1 electron--vibron coupling
$M = D_0\,\chi(z)\,|u\rangle\langle u|$.
(d)~Schematic of the rank-1 projected SCBA loop: a single
scalar-in-mode-space, 1D-in-$z$ Dyson equation dressed self-consistently by
the phonon self-energies $\tilde\sigma^{R,<,>}$. This replaces the full
nine-mode coupled SCBA with machine-precision accuracy whenever
$\sigma_{\mathrm{mol}} k_\perp \ll 1$.}
\label{fig:device}
\end{figure}

% --- Fig. 2 ----------------------------------------------------------------
\begin{figure}[tbp]
\centerline{\includegraphics[width=0.95\columnwidth]{figs/fig2_patil.pdf}}
\caption{\textbf{Validation against the Patil 1D benchmark (Tier-1 $\gamma$,
$\delta$).}
Right-contact current $I_R(V)$ of Patil's 1D symmetric RTD computed with
the rank-1 Keldysh wrapper (red circles) overlaid on the Patil 1D reference
from the published MATLAB code \texttt{rtd2modes\_1d.m} (black line). The
rank-1 solver reproduces the NDR peak to 0.55\%:
$I_R = 5.667 \times 10^{-10}$~A (rank-1) versus
$5.698 \times 10^{-10}$~A (reference) at $V = 0.336$~V. The orange
triangles show $-I_L$ on the same axes; the two traces would be exactly
coincident for a perfectly converged SCBA solution.
Inset: the per-bias current-conservation residual
$|I_L + I_R| / \max(|I_L|,|I_R|)$ in log scale. Exposing both contacts lets
us read $\delta$ directly from the Meir--Wingreen trace without any
auxiliary bookkeeping. All parameters (device, mesh, bias convention,
phonon modes, coupling, temperature) are identical to the reference.}
\label{fig:patil}
\end{figure}

% --- Fig. 3 ----------------------------------------------------------------
\begin{figure}[tbp]
\centerline{\includegraphics[width=0.95\columnwidth]{figs/fig3_iv.pdf}}
\caption{\textbf{Current--voltage characteristics of the ZnO/MgZnO primary
stack at 300~K.}
Rank-1 projected SCBA sweep (\texttt{max\_iter}$=10$, mix$=0.4$) on the
asymmetric ZnO/Mg\textsubscript{0.3}Zn\textsubscript{0.7}O RTD for the
four test species: Baseline (blank reference, bulk ZnO LO phonon at 72~meV
only, no molecule), Mol\_A (single vibrational mode $\hbar\omega_A =
100$~meV), Mol\_B ($\hbar\omega_B = 180$~meV), and Mol\_AB (both modes).
All molecules couple with $D_0^2 = 0.1\,\mathrm{eV}^2$ through a Gaussian
form factor of width 3~\AA\ centered on the emitter-barrier midpoint. Each
trace shows a clean NDR feature at $V \approx V_{\mathrm{res}}$; the phonon
sidebands pull current into the sub-resonance region, shifting the turn-on
and broadening the NDR step in a mode-specific way. Bias grid: 201~points
over 0--400~mV. Energy grid: $dE = 2$~meV over $[-0.25, +0.5]$~eV.
$E_F = 20$~meV. No numerical differentiation is applied to $I(V)$
anywhere in this work; the $d^2I/dV^2$ in Fig.~\ref{fig:d2i} comes from
the closed-form formula~(\ref{eq:d2i}) evaluated on the same converged
Keldysh state used to build this figure.}
\label{fig:iv}
\end{figure}

% --- Fig. 4 ----------------------------------------------------------------
\begin{figure}[tbp]
\centerline{\includegraphics[width=0.95\columnwidth]{figs/fig4_d2i.pdf}}
\caption{\textbf{Analytic $d^2I_R/dV^2$ --- vibrational fingerprints at
300~K.}
The second derivative is computed in closed form directly from the
converged Keldysh state at each bias via
$d^2I_R/dV^2 = (e^2/h)\cdot(1/4)\int dE\,[f_L''(E,V)\,T_{RL}(E) +
f_R''(E,V)\,(T_{RR}(E) - T_{RA}(E))]$, with the kernels
$T_{RL}, T_{RR}, T_{RA}$ built once from
$(G^R,\Gamma_L,\Gamma_R)$ at each bias --- no numerical $V$-grid
differentiation. At $D_0=0$ the formula reduces to the Landauer coherent
expression by the optical theorem
$T_{RA} = T_{RL} + T_{RR}$ (unit-tested to $10^{-6}$).
Mol\_A shows a peak at $V \approx V_{\mathrm{res}} + \hbar\omega_A/e \approx
187$~mV; Mol\_B at $V \approx V_{\mathrm{res}} + \hbar\omega_B/e \approx
267$~mV; Mol\_AB shows both features simultaneously. Baseline has a shared
feature at $\sim 162$~mV from the 72~meV bulk LO phonon. Red ticks mark the
predicted sideband biases: these are \emph{not fits}.}
\label{fig:d2i}
\end{figure}

% --- Fig. 5 ----------------------------------------------------------------
\begin{figure}[tbp]
\centerline{\includegraphics[width=0.95\columnwidth]{figs/fig5_discrimination.pdf}}
\caption{\textbf{Discrimination metric
$\Delta D(V) = d^2I_{\mathrm{mol}}/dV^2 - d^2I_{\mathrm{Baseline}}/dV^2$.}
Baseline-subtracted second-derivative spectra isolate the molecular IETS
contribution from the bulk-phonon and coherent backgrounds. The 100~meV
and 180~meV modes produce non-overlapping features separated by 80~mV,
well outside the thermal broadening $3k_BT/e \approx 78$~mV at 300~K. For
Mol\_AB the two features add linearly ($\Delta D_{AB}(V) \approx
\Delta D_A(V) + \Delta D_B(V)$ within numerical noise), a direct
demonstration of the closed-form analytic second-derivative (\ref{eq:d2i})
acting on a mixed analyte. The peak-to-peak values
$|\Delta D|_{\max}$ for each single-mode species are the figures of merit
for pixel-level discrimination. No lock-in or smoothing of any kind is
applied: the signal here is the direct output of the closed-form
(\ref{eq:d2i}) minus its Baseline evaluation.}
\label{fig:disc}
\end{figure}

%% or drop the individual \begin{figure}…\end{figure} blocks inline.
%% ----------------------------------------------------------------------

% --- Fig. 1 ----------------------------------------------------------------
\begin{figure}[tbp]
\centerline{\includegraphics[width=0.95\columnwidth]{figs/fig1_device.pdf}}
\caption{\textbf{Device and methodology at a glance.}
(a)~Layer stack of the asymmetric
ZnO/Mg\textsubscript{0.3}Zn\textsubscript{0.7}O RTD: 10~nm $n^+$-ZnO emitter
contact, 3~nm Mg\textsubscript{0.3}Zn\textsubscript{0.7}O emitter barrier
(host of the sensing molecule), 3~nm ZnO quantum well, 1.5~nm
Mg\textsubscript{0.3}Zn\textsubscript{0.7}O collector barrier, and 10~nm
$n^+$-ZnO collector. Conduction-band offset $\Delta E_c = 0.47$~eV,
effective mass $m^* = 0.28\,m_0$. The transverse area is
$A = 10\,\mu\mathrm{m} \times 10\,\mu\mathrm{m}$, entering the current
expression as a Datta ``Fix-$A$'' prefactor.
(b)~Conduction-band profile at $V = 0$ and at $V = V_{\mathrm{res}}$, showing
the quantized well level $E_{\mathrm{res}}$ crossing $E_F$ from above to
give negative differential resistance.
(c)~Molecular form factor $\chi(z)$, a Gaussian of width
$\sigma_{\mathrm{mol}} = 3$~\AA\ centered at the emitter-barrier midpoint,
defining the rank-1 electron--vibron coupling
$M = D_0\,\chi(z)\,|u\rangle\langle u|$.
(d)~Schematic of the rank-1 projected SCBA loop: a single
scalar-in-mode-space, 1D-in-$z$ Dyson equation dressed self-consistently by
the phonon self-energies $\tilde\sigma^{R,<,>}$. This replaces the full
nine-mode coupled SCBA with machine-precision accuracy whenever
$\sigma_{\mathrm{mol}} k_\perp \ll 1$.}
\label{fig:device}
\end{figure}

% --- Fig. 2 ----------------------------------------------------------------
\begin{figure}[tbp]
\centerline{\includegraphics[width=0.95\columnwidth]{figs/fig2_patil.pdf}}
\caption{\textbf{Validation against the Patil 1D benchmark (Tier-1 $\gamma$,
$\delta$).}
Right-contact current $I_R(V)$ of Patil's 1D symmetric RTD computed with
the rank-1 Keldysh wrapper (red circles) overlaid on the Patil 1D reference
from the published MATLAB code \texttt{rtd2modes\_1d.m} (black line). The
rank-1 solver reproduces the NDR peak to 0.55\%:
$I_R = 5.667 \times 10^{-10}$~A (rank-1) versus
$5.698 \times 10^{-10}$~A (reference) at $V = 0.336$~V. The orange
triangles show $-I_L$ on the same axes; the two traces would be exactly
coincident for a perfectly converged SCBA solution.
Inset: the per-bias current-conservation residual
$|I_L + I_R| / \max(|I_L|,|I_R|)$ in log scale. Exposing both contacts lets
us read $\delta$ directly from the Meir--Wingreen trace without any
auxiliary bookkeeping. All parameters (device, mesh, bias convention,
phonon modes, coupling, temperature) are identical to the reference.}
\label{fig:patil}
\end{figure}

% --- Fig. 3 ----------------------------------------------------------------
\begin{figure}[tbp]
\centerline{\includegraphics[width=0.95\columnwidth]{figs/fig3_iv.pdf}}
\caption{\textbf{Current--voltage characteristics of the ZnO/MgZnO primary
stack at 300~K.}
Rank-1 projected SCBA sweep (\texttt{max\_iter}$=10$, mix$=0.4$) on the
asymmetric ZnO/Mg\textsubscript{0.3}Zn\textsubscript{0.7}O RTD for the
four test species: Baseline (blank reference, bulk ZnO LO phonon at 72~meV
only, no molecule), Mol\_A (single vibrational mode $\hbar\omega_A =
100$~meV), Mol\_B ($\hbar\omega_B = 180$~meV), and Mol\_AB (both modes).
All molecules couple with $D_0^2 = 0.1\,\mathrm{eV}^2$ through a Gaussian
form factor of width 3~\AA\ centered on the emitter-barrier midpoint. Each
trace shows a clean NDR feature at $V \approx V_{\mathrm{res}}$; the phonon
sidebands pull current into the sub-resonance region, shifting the turn-on
and broadening the NDR step in a mode-specific way. Bias grid: 201~points
over 0--400~mV. Energy grid: $dE = 2$~meV over $[-0.25, +0.5]$~eV.
$E_F = 20$~meV. No numerical differentiation is applied to $I(V)$
anywhere in this work; the $d^2I/dV^2$ in Fig.~\ref{fig:d2i} comes from
the closed-form formula~(\ref{eq:d2i}) evaluated on the same converged
Keldysh state used to build this figure.}
\label{fig:iv}
\end{figure}

% --- Fig. 4 ----------------------------------------------------------------
\begin{figure}[tbp]
\centerline{\includegraphics[width=0.95\columnwidth]{figs/fig4_d2i.pdf}}
\caption{\textbf{Analytic $d^2I_R/dV^2$ --- vibrational fingerprints at
300~K.}
The second derivative is computed in closed form directly from the
converged Keldysh state at each bias via
$d^2I_R/dV^2 = (e^2/h)\cdot(1/4)\int dE\,[f_L''(E,V)\,T_{RL}(E) +
f_R''(E,V)\,(T_{RR}(E) - T_{RA}(E))]$, with the kernels
$T_{RL}, T_{RR}, T_{RA}$ built once from
$(G^R,\Gamma_L,\Gamma_R)$ at each bias --- no numerical $V$-grid
differentiation. At $D_0=0$ the formula reduces to the Landauer coherent
expression by the optical theorem
$T_{RA} = T_{RL} + T_{RR}$ (unit-tested to $10^{-6}$).
Mol\_A shows a peak at $V \approx V_{\mathrm{res}} + \hbar\omega_A/e \approx
187$~mV; Mol\_B at $V \approx V_{\mathrm{res}} + \hbar\omega_B/e \approx
267$~mV; Mol\_AB shows both features simultaneously. Baseline has a shared
feature at $\sim 162$~mV from the 72~meV bulk LO phonon. Red ticks mark the
predicted sideband biases: these are \emph{not fits}.}
\label{fig:d2i}
\end{figure}

% --- Fig. 5 ----------------------------------------------------------------
\begin{figure}[tbp]
\centerline{\includegraphics[width=0.95\columnwidth]{figs/fig5_discrimination.pdf}}
\caption{\textbf{Discrimination metric
$\Delta D(V) = d^2I_{\mathrm{mol}}/dV^2 - d^2I_{\mathrm{Baseline}}/dV^2$.}
Baseline-subtracted second-derivative spectra isolate the molecular IETS
contribution from the bulk-phonon and coherent backgrounds. The 100~meV
and 180~meV modes produce non-overlapping features separated by 80~mV,
well outside the thermal broadening $3k_BT/e \approx 78$~mV at 300~K. For
Mol\_AB the two features add linearly ($\Delta D_{AB}(V) \approx
\Delta D_A(V) + \Delta D_B(V)$ within numerical noise), a direct
demonstration of the closed-form analytic second-derivative (\ref{eq:d2i})
acting on a mixed analyte. The peak-to-peak values
$|\Delta D|_{\max}$ for each single-mode species are the figures of merit
for pixel-level discrimination. No lock-in or smoothing of any kind is
applied: the signal here is the direct output of the closed-form
(\ref{eq:d2i}) minus its Baseline evaluation.}
\label{fig:disc}
\end{figure}
.
%%
%% Each caption is written to be fully self-contained: an expert in the
%% field who reads only the paper title and this figure + caption should
%% be able to tell what was done, how, and what was shown.
%%
%% Figure files (to be generated from results/2026-04-09/*.npz by
%% run/plot_sispad_figs.py):
%%
%%   figs/fig1_device.pdf     — Device stack, band diagram, chi(z), rank-1 loop
%%   figs/fig2_patil.pdf      — Patil 1D benchmark + conservation residual inset
%%   figs/fig3_iv.pdf         — I–V of Baseline/Mol_A/Mol_B/Mol_AB on ZnO/MgZnO
%%   figs/fig4_d2i.pdf        — Analytic d2I/dV2 for all 4 molecules
%%   figs/fig5_discrimination.pdf — Baseline-subtracted discrimination metric
%%
%% Usage in sispad2026_abstract.tex:
%%   %% SISPAD 2026 — figures + captions, kept separate from the main abstract
%% file so captions can be iterated independently and dropped into
%% papers/sispad2026_abstract.tex via %% SISPAD 2026 — figures + captions, kept separate from the main abstract
%% file so captions can be iterated independently and dropped into
%% papers/sispad2026_abstract.tex via %% SISPAD 2026 — figures + captions, kept separate from the main abstract
%% file so captions can be iterated independently and dropped into
%% papers/sispad2026_abstract.tex via \input{sispad2026_figures.tex}.
%%
%% Each caption is written to be fully self-contained: an expert in the
%% field who reads only the paper title and this figure + caption should
%% be able to tell what was done, how, and what was shown.
%%
%% Figure files (to be generated from results/2026-04-09/*.npz by
%% run/plot_sispad_figs.py):
%%
%%   figs/fig1_device.pdf     — Device stack, band diagram, chi(z), rank-1 loop
%%   figs/fig2_patil.pdf      — Patil 1D benchmark + conservation residual inset
%%   figs/fig3_iv.pdf         — I–V of Baseline/Mol_A/Mol_B/Mol_AB on ZnO/MgZnO
%%   figs/fig4_d2i.pdf        — Analytic d2I/dV2 for all 4 molecules
%%   figs/fig5_discrimination.pdf — Baseline-subtracted discrimination metric
%%
%% Usage in sispad2026_abstract.tex:
%%   \input{sispad2026_figures}
%% or drop the individual \begin{figure}…\end{figure} blocks inline.
%% ----------------------------------------------------------------------

% --- Fig. 1 ----------------------------------------------------------------
\begin{figure}[tbp]
\centerline{\includegraphics[width=0.95\columnwidth]{figs/fig1_device.pdf}}
\caption{\textbf{Device and methodology at a glance.}
(a)~Layer stack of the asymmetric
ZnO/Mg\textsubscript{0.3}Zn\textsubscript{0.7}O RTD: 10~nm $n^+$-ZnO emitter
contact, 3~nm Mg\textsubscript{0.3}Zn\textsubscript{0.7}O emitter barrier
(host of the sensing molecule), 3~nm ZnO quantum well, 1.5~nm
Mg\textsubscript{0.3}Zn\textsubscript{0.7}O collector barrier, and 10~nm
$n^+$-ZnO collector. Conduction-band offset $\Delta E_c = 0.47$~eV,
effective mass $m^* = 0.28\,m_0$. The transverse area is
$A = 10\,\mu\mathrm{m} \times 10\,\mu\mathrm{m}$, entering the current
expression as a Datta ``Fix-$A$'' prefactor.
(b)~Conduction-band profile at $V = 0$ and at $V = V_{\mathrm{res}}$, showing
the quantized well level $E_{\mathrm{res}}$ crossing $E_F$ from above to
give negative differential resistance.
(c)~Molecular form factor $\chi(z)$, a Gaussian of width
$\sigma_{\mathrm{mol}} = 3$~\AA\ centered at the emitter-barrier midpoint,
defining the rank-1 electron--vibron coupling
$M = D_0\,\chi(z)\,|u\rangle\langle u|$.
(d)~Schematic of the rank-1 projected SCBA loop: a single
scalar-in-mode-space, 1D-in-$z$ Dyson equation dressed self-consistently by
the phonon self-energies $\tilde\sigma^{R,<,>}$. This replaces the full
nine-mode coupled SCBA with machine-precision accuracy whenever
$\sigma_{\mathrm{mol}} k_\perp \ll 1$.}
\label{fig:device}
\end{figure}

% --- Fig. 2 ----------------------------------------------------------------
\begin{figure}[tbp]
\centerline{\includegraphics[width=0.95\columnwidth]{figs/fig2_patil.pdf}}
\caption{\textbf{Validation against the Patil 1D benchmark (Tier-1 $\gamma$,
$\delta$).}
Right-contact current $I_R(V)$ of Patil's 1D symmetric RTD computed with
the rank-1 Keldysh wrapper (red circles) overlaid on the Patil 1D reference
from the published MATLAB code \texttt{rtd2modes\_1d.m} (black line). The
rank-1 solver reproduces the NDR peak to 0.55\%:
$I_R = 5.667 \times 10^{-10}$~A (rank-1) versus
$5.698 \times 10^{-10}$~A (reference) at $V = 0.336$~V. The orange
triangles show $-I_L$ on the same axes; the two traces would be exactly
coincident for a perfectly converged SCBA solution.
Inset: the per-bias current-conservation residual
$|I_L + I_R| / \max(|I_L|,|I_R|)$ in log scale. Exposing both contacts lets
us read $\delta$ directly from the Meir--Wingreen trace without any
auxiliary bookkeeping. All parameters (device, mesh, bias convention,
phonon modes, coupling, temperature) are identical to the reference.}
\label{fig:patil}
\end{figure}

% --- Fig. 3 ----------------------------------------------------------------
\begin{figure}[tbp]
\centerline{\includegraphics[width=0.95\columnwidth]{figs/fig3_iv.pdf}}
\caption{\textbf{Current--voltage characteristics of the ZnO/MgZnO primary
stack at 300~K.}
Rank-1 projected SCBA sweep (\texttt{max\_iter}$=10$, mix$=0.4$) on the
asymmetric ZnO/Mg\textsubscript{0.3}Zn\textsubscript{0.7}O RTD for the
four test species: Baseline (blank reference, bulk ZnO LO phonon at 72~meV
only, no molecule), Mol\_A (single vibrational mode $\hbar\omega_A =
100$~meV), Mol\_B ($\hbar\omega_B = 180$~meV), and Mol\_AB (both modes).
All molecules couple with $D_0^2 = 0.1\,\mathrm{eV}^2$ through a Gaussian
form factor of width 3~\AA\ centered on the emitter-barrier midpoint. Each
trace shows a clean NDR feature at $V \approx V_{\mathrm{res}}$; the phonon
sidebands pull current into the sub-resonance region, shifting the turn-on
and broadening the NDR step in a mode-specific way. Bias grid: 201~points
over 0--400~mV. Energy grid: $dE = 2$~meV over $[-0.25, +0.5]$~eV.
$E_F = 20$~meV. No numerical differentiation is applied to $I(V)$
anywhere in this work; the $d^2I/dV^2$ in Fig.~\ref{fig:d2i} comes from
the closed-form formula~(\ref{eq:d2i}) evaluated on the same converged
Keldysh state used to build this figure.}
\label{fig:iv}
\end{figure}

% --- Fig. 4 ----------------------------------------------------------------
\begin{figure}[tbp]
\centerline{\includegraphics[width=0.95\columnwidth]{figs/fig4_d2i.pdf}}
\caption{\textbf{Analytic $d^2I_R/dV^2$ --- vibrational fingerprints at
300~K.}
The second derivative is computed in closed form directly from the
converged Keldysh state at each bias via
$d^2I_R/dV^2 = (e^2/h)\cdot(1/4)\int dE\,[f_L''(E,V)\,T_{RL}(E) +
f_R''(E,V)\,(T_{RR}(E) - T_{RA}(E))]$, with the kernels
$T_{RL}, T_{RR}, T_{RA}$ built once from
$(G^R,\Gamma_L,\Gamma_R)$ at each bias --- no numerical $V$-grid
differentiation. At $D_0=0$ the formula reduces to the Landauer coherent
expression by the optical theorem
$T_{RA} = T_{RL} + T_{RR}$ (unit-tested to $10^{-6}$).
Mol\_A shows a peak at $V \approx V_{\mathrm{res}} + \hbar\omega_A/e \approx
187$~mV; Mol\_B at $V \approx V_{\mathrm{res}} + \hbar\omega_B/e \approx
267$~mV; Mol\_AB shows both features simultaneously. Baseline has a shared
feature at $\sim 162$~mV from the 72~meV bulk LO phonon. Red ticks mark the
predicted sideband biases: these are \emph{not fits}.}
\label{fig:d2i}
\end{figure}

% --- Fig. 5 ----------------------------------------------------------------
\begin{figure}[tbp]
\centerline{\includegraphics[width=0.95\columnwidth]{figs/fig5_discrimination.pdf}}
\caption{\textbf{Discrimination metric
$\Delta D(V) = d^2I_{\mathrm{mol}}/dV^2 - d^2I_{\mathrm{Baseline}}/dV^2$.}
Baseline-subtracted second-derivative spectra isolate the molecular IETS
contribution from the bulk-phonon and coherent backgrounds. The 100~meV
and 180~meV modes produce non-overlapping features separated by 80~mV,
well outside the thermal broadening $3k_BT/e \approx 78$~mV at 300~K. For
Mol\_AB the two features add linearly ($\Delta D_{AB}(V) \approx
\Delta D_A(V) + \Delta D_B(V)$ within numerical noise), a direct
demonstration of the closed-form analytic second-derivative (\ref{eq:d2i})
acting on a mixed analyte. The peak-to-peak values
$|\Delta D|_{\max}$ for each single-mode species are the figures of merit
for pixel-level discrimination. No lock-in or smoothing of any kind is
applied: the signal here is the direct output of the closed-form
(\ref{eq:d2i}) minus its Baseline evaluation.}
\label{fig:disc}
\end{figure}
.
%%
%% Each caption is written to be fully self-contained: an expert in the
%% field who reads only the paper title and this figure + caption should
%% be able to tell what was done, how, and what was shown.
%%
%% Figure files (to be generated from results/2026-04-09/*.npz by
%% run/plot_sispad_figs.py):
%%
%%   figs/fig1_device.pdf     — Device stack, band diagram, chi(z), rank-1 loop
%%   figs/fig2_patil.pdf      — Patil 1D benchmark + conservation residual inset
%%   figs/fig3_iv.pdf         — I–V of Baseline/Mol_A/Mol_B/Mol_AB on ZnO/MgZnO
%%   figs/fig4_d2i.pdf        — Analytic d2I/dV2 for all 4 molecules
%%   figs/fig5_discrimination.pdf — Baseline-subtracted discrimination metric
%%
%% Usage in sispad2026_abstract.tex:
%%   %% SISPAD 2026 — figures + captions, kept separate from the main abstract
%% file so captions can be iterated independently and dropped into
%% papers/sispad2026_abstract.tex via \input{sispad2026_figures.tex}.
%%
%% Each caption is written to be fully self-contained: an expert in the
%% field who reads only the paper title and this figure + caption should
%% be able to tell what was done, how, and what was shown.
%%
%% Figure files (to be generated from results/2026-04-09/*.npz by
%% run/plot_sispad_figs.py):
%%
%%   figs/fig1_device.pdf     — Device stack, band diagram, chi(z), rank-1 loop
%%   figs/fig2_patil.pdf      — Patil 1D benchmark + conservation residual inset
%%   figs/fig3_iv.pdf         — I–V of Baseline/Mol_A/Mol_B/Mol_AB on ZnO/MgZnO
%%   figs/fig4_d2i.pdf        — Analytic d2I/dV2 for all 4 molecules
%%   figs/fig5_discrimination.pdf — Baseline-subtracted discrimination metric
%%
%% Usage in sispad2026_abstract.tex:
%%   \input{sispad2026_figures}
%% or drop the individual \begin{figure}…\end{figure} blocks inline.
%% ----------------------------------------------------------------------

% --- Fig. 1 ----------------------------------------------------------------
\begin{figure}[tbp]
\centerline{\includegraphics[width=0.95\columnwidth]{figs/fig1_device.pdf}}
\caption{\textbf{Device and methodology at a glance.}
(a)~Layer stack of the asymmetric
ZnO/Mg\textsubscript{0.3}Zn\textsubscript{0.7}O RTD: 10~nm $n^+$-ZnO emitter
contact, 3~nm Mg\textsubscript{0.3}Zn\textsubscript{0.7}O emitter barrier
(host of the sensing molecule), 3~nm ZnO quantum well, 1.5~nm
Mg\textsubscript{0.3}Zn\textsubscript{0.7}O collector barrier, and 10~nm
$n^+$-ZnO collector. Conduction-band offset $\Delta E_c = 0.47$~eV,
effective mass $m^* = 0.28\,m_0$. The transverse area is
$A = 10\,\mu\mathrm{m} \times 10\,\mu\mathrm{m}$, entering the current
expression as a Datta ``Fix-$A$'' prefactor.
(b)~Conduction-band profile at $V = 0$ and at $V = V_{\mathrm{res}}$, showing
the quantized well level $E_{\mathrm{res}}$ crossing $E_F$ from above to
give negative differential resistance.
(c)~Molecular form factor $\chi(z)$, a Gaussian of width
$\sigma_{\mathrm{mol}} = 3$~\AA\ centered at the emitter-barrier midpoint,
defining the rank-1 electron--vibron coupling
$M = D_0\,\chi(z)\,|u\rangle\langle u|$.
(d)~Schematic of the rank-1 projected SCBA loop: a single
scalar-in-mode-space, 1D-in-$z$ Dyson equation dressed self-consistently by
the phonon self-energies $\tilde\sigma^{R,<,>}$. This replaces the full
nine-mode coupled SCBA with machine-precision accuracy whenever
$\sigma_{\mathrm{mol}} k_\perp \ll 1$.}
\label{fig:device}
\end{figure}

% --- Fig. 2 ----------------------------------------------------------------
\begin{figure}[tbp]
\centerline{\includegraphics[width=0.95\columnwidth]{figs/fig2_patil.pdf}}
\caption{\textbf{Validation against the Patil 1D benchmark (Tier-1 $\gamma$,
$\delta$).}
Right-contact current $I_R(V)$ of Patil's 1D symmetric RTD computed with
the rank-1 Keldysh wrapper (red circles) overlaid on the Patil 1D reference
from the published MATLAB code \texttt{rtd2modes\_1d.m} (black line). The
rank-1 solver reproduces the NDR peak to 0.55\%:
$I_R = 5.667 \times 10^{-10}$~A (rank-1) versus
$5.698 \times 10^{-10}$~A (reference) at $V = 0.336$~V. The orange
triangles show $-I_L$ on the same axes; the two traces would be exactly
coincident for a perfectly converged SCBA solution.
Inset: the per-bias current-conservation residual
$|I_L + I_R| / \max(|I_L|,|I_R|)$ in log scale. Exposing both contacts lets
us read $\delta$ directly from the Meir--Wingreen trace without any
auxiliary bookkeeping. All parameters (device, mesh, bias convention,
phonon modes, coupling, temperature) are identical to the reference.}
\label{fig:patil}
\end{figure}

% --- Fig. 3 ----------------------------------------------------------------
\begin{figure}[tbp]
\centerline{\includegraphics[width=0.95\columnwidth]{figs/fig3_iv.pdf}}
\caption{\textbf{Current--voltage characteristics of the ZnO/MgZnO primary
stack at 300~K.}
Rank-1 projected SCBA sweep (\texttt{max\_iter}$=10$, mix$=0.4$) on the
asymmetric ZnO/Mg\textsubscript{0.3}Zn\textsubscript{0.7}O RTD for the
four test species: Baseline (blank reference, bulk ZnO LO phonon at 72~meV
only, no molecule), Mol\_A (single vibrational mode $\hbar\omega_A =
100$~meV), Mol\_B ($\hbar\omega_B = 180$~meV), and Mol\_AB (both modes).
All molecules couple with $D_0^2 = 0.1\,\mathrm{eV}^2$ through a Gaussian
form factor of width 3~\AA\ centered on the emitter-barrier midpoint. Each
trace shows a clean NDR feature at $V \approx V_{\mathrm{res}}$; the phonon
sidebands pull current into the sub-resonance region, shifting the turn-on
and broadening the NDR step in a mode-specific way. Bias grid: 201~points
over 0--400~mV. Energy grid: $dE = 2$~meV over $[-0.25, +0.5]$~eV.
$E_F = 20$~meV. No numerical differentiation is applied to $I(V)$
anywhere in this work; the $d^2I/dV^2$ in Fig.~\ref{fig:d2i} comes from
the closed-form formula~(\ref{eq:d2i}) evaluated on the same converged
Keldysh state used to build this figure.}
\label{fig:iv}
\end{figure}

% --- Fig. 4 ----------------------------------------------------------------
\begin{figure}[tbp]
\centerline{\includegraphics[width=0.95\columnwidth]{figs/fig4_d2i.pdf}}
\caption{\textbf{Analytic $d^2I_R/dV^2$ --- vibrational fingerprints at
300~K.}
The second derivative is computed in closed form directly from the
converged Keldysh state at each bias via
$d^2I_R/dV^2 = (e^2/h)\cdot(1/4)\int dE\,[f_L''(E,V)\,T_{RL}(E) +
f_R''(E,V)\,(T_{RR}(E) - T_{RA}(E))]$, with the kernels
$T_{RL}, T_{RR}, T_{RA}$ built once from
$(G^R,\Gamma_L,\Gamma_R)$ at each bias --- no numerical $V$-grid
differentiation. At $D_0=0$ the formula reduces to the Landauer coherent
expression by the optical theorem
$T_{RA} = T_{RL} + T_{RR}$ (unit-tested to $10^{-6}$).
Mol\_A shows a peak at $V \approx V_{\mathrm{res}} + \hbar\omega_A/e \approx
187$~mV; Mol\_B at $V \approx V_{\mathrm{res}} + \hbar\omega_B/e \approx
267$~mV; Mol\_AB shows both features simultaneously. Baseline has a shared
feature at $\sim 162$~mV from the 72~meV bulk LO phonon. Red ticks mark the
predicted sideband biases: these are \emph{not fits}.}
\label{fig:d2i}
\end{figure}

% --- Fig. 5 ----------------------------------------------------------------
\begin{figure}[tbp]
\centerline{\includegraphics[width=0.95\columnwidth]{figs/fig5_discrimination.pdf}}
\caption{\textbf{Discrimination metric
$\Delta D(V) = d^2I_{\mathrm{mol}}/dV^2 - d^2I_{\mathrm{Baseline}}/dV^2$.}
Baseline-subtracted second-derivative spectra isolate the molecular IETS
contribution from the bulk-phonon and coherent backgrounds. The 100~meV
and 180~meV modes produce non-overlapping features separated by 80~mV,
well outside the thermal broadening $3k_BT/e \approx 78$~mV at 300~K. For
Mol\_AB the two features add linearly ($\Delta D_{AB}(V) \approx
\Delta D_A(V) + \Delta D_B(V)$ within numerical noise), a direct
demonstration of the closed-form analytic second-derivative (\ref{eq:d2i})
acting on a mixed analyte. The peak-to-peak values
$|\Delta D|_{\max}$ for each single-mode species are the figures of merit
for pixel-level discrimination. No lock-in or smoothing of any kind is
applied: the signal here is the direct output of the closed-form
(\ref{eq:d2i}) minus its Baseline evaluation.}
\label{fig:disc}
\end{figure}

%% or drop the individual \begin{figure}…\end{figure} blocks inline.
%% ----------------------------------------------------------------------

% --- Fig. 1 ----------------------------------------------------------------
\begin{figure}[tbp]
\centerline{\includegraphics[width=0.95\columnwidth]{figs/fig1_device.pdf}}
\caption{\textbf{Device and methodology at a glance.}
(a)~Layer stack of the asymmetric
ZnO/Mg\textsubscript{0.3}Zn\textsubscript{0.7}O RTD: 10~nm $n^+$-ZnO emitter
contact, 3~nm Mg\textsubscript{0.3}Zn\textsubscript{0.7}O emitter barrier
(host of the sensing molecule), 3~nm ZnO quantum well, 1.5~nm
Mg\textsubscript{0.3}Zn\textsubscript{0.7}O collector barrier, and 10~nm
$n^+$-ZnO collector. Conduction-band offset $\Delta E_c = 0.47$~eV,
effective mass $m^* = 0.28\,m_0$. The transverse area is
$A = 10\,\mu\mathrm{m} \times 10\,\mu\mathrm{m}$, entering the current
expression as a Datta ``Fix-$A$'' prefactor.
(b)~Conduction-band profile at $V = 0$ and at $V = V_{\mathrm{res}}$, showing
the quantized well level $E_{\mathrm{res}}$ crossing $E_F$ from above to
give negative differential resistance.
(c)~Molecular form factor $\chi(z)$, a Gaussian of width
$\sigma_{\mathrm{mol}} = 3$~\AA\ centered at the emitter-barrier midpoint,
defining the rank-1 electron--vibron coupling
$M = D_0\,\chi(z)\,|u\rangle\langle u|$.
(d)~Schematic of the rank-1 projected SCBA loop: a single
scalar-in-mode-space, 1D-in-$z$ Dyson equation dressed self-consistently by
the phonon self-energies $\tilde\sigma^{R,<,>}$. This replaces the full
nine-mode coupled SCBA with machine-precision accuracy whenever
$\sigma_{\mathrm{mol}} k_\perp \ll 1$.}
\label{fig:device}
\end{figure}

% --- Fig. 2 ----------------------------------------------------------------
\begin{figure}[tbp]
\centerline{\includegraphics[width=0.95\columnwidth]{figs/fig2_patil.pdf}}
\caption{\textbf{Validation against the Patil 1D benchmark (Tier-1 $\gamma$,
$\delta$).}
Right-contact current $I_R(V)$ of Patil's 1D symmetric RTD computed with
the rank-1 Keldysh wrapper (red circles) overlaid on the Patil 1D reference
from the published MATLAB code \texttt{rtd2modes\_1d.m} (black line). The
rank-1 solver reproduces the NDR peak to 0.55\%:
$I_R = 5.667 \times 10^{-10}$~A (rank-1) versus
$5.698 \times 10^{-10}$~A (reference) at $V = 0.336$~V. The orange
triangles show $-I_L$ on the same axes; the two traces would be exactly
coincident for a perfectly converged SCBA solution.
Inset: the per-bias current-conservation residual
$|I_L + I_R| / \max(|I_L|,|I_R|)$ in log scale. Exposing both contacts lets
us read $\delta$ directly from the Meir--Wingreen trace without any
auxiliary bookkeeping. All parameters (device, mesh, bias convention,
phonon modes, coupling, temperature) are identical to the reference.}
\label{fig:patil}
\end{figure}

% --- Fig. 3 ----------------------------------------------------------------
\begin{figure}[tbp]
\centerline{\includegraphics[width=0.95\columnwidth]{figs/fig3_iv.pdf}}
\caption{\textbf{Current--voltage characteristics of the ZnO/MgZnO primary
stack at 300~K.}
Rank-1 projected SCBA sweep (\texttt{max\_iter}$=10$, mix$=0.4$) on the
asymmetric ZnO/Mg\textsubscript{0.3}Zn\textsubscript{0.7}O RTD for the
four test species: Baseline (blank reference, bulk ZnO LO phonon at 72~meV
only, no molecule), Mol\_A (single vibrational mode $\hbar\omega_A =
100$~meV), Mol\_B ($\hbar\omega_B = 180$~meV), and Mol\_AB (both modes).
All molecules couple with $D_0^2 = 0.1\,\mathrm{eV}^2$ through a Gaussian
form factor of width 3~\AA\ centered on the emitter-barrier midpoint. Each
trace shows a clean NDR feature at $V \approx V_{\mathrm{res}}$; the phonon
sidebands pull current into the sub-resonance region, shifting the turn-on
and broadening the NDR step in a mode-specific way. Bias grid: 201~points
over 0--400~mV. Energy grid: $dE = 2$~meV over $[-0.25, +0.5]$~eV.
$E_F = 20$~meV. No numerical differentiation is applied to $I(V)$
anywhere in this work; the $d^2I/dV^2$ in Fig.~\ref{fig:d2i} comes from
the closed-form formula~(\ref{eq:d2i}) evaluated on the same converged
Keldysh state used to build this figure.}
\label{fig:iv}
\end{figure}

% --- Fig. 4 ----------------------------------------------------------------
\begin{figure}[tbp]
\centerline{\includegraphics[width=0.95\columnwidth]{figs/fig4_d2i.pdf}}
\caption{\textbf{Analytic $d^2I_R/dV^2$ --- vibrational fingerprints at
300~K.}
The second derivative is computed in closed form directly from the
converged Keldysh state at each bias via
$d^2I_R/dV^2 = (e^2/h)\cdot(1/4)\int dE\,[f_L''(E,V)\,T_{RL}(E) +
f_R''(E,V)\,(T_{RR}(E) - T_{RA}(E))]$, with the kernels
$T_{RL}, T_{RR}, T_{RA}$ built once from
$(G^R,\Gamma_L,\Gamma_R)$ at each bias --- no numerical $V$-grid
differentiation. At $D_0=0$ the formula reduces to the Landauer coherent
expression by the optical theorem
$T_{RA} = T_{RL} + T_{RR}$ (unit-tested to $10^{-6}$).
Mol\_A shows a peak at $V \approx V_{\mathrm{res}} + \hbar\omega_A/e \approx
187$~mV; Mol\_B at $V \approx V_{\mathrm{res}} + \hbar\omega_B/e \approx
267$~mV; Mol\_AB shows both features simultaneously. Baseline has a shared
feature at $\sim 162$~mV from the 72~meV bulk LO phonon. Red ticks mark the
predicted sideband biases: these are \emph{not fits}.}
\label{fig:d2i}
\end{figure}

% --- Fig. 5 ----------------------------------------------------------------
\begin{figure}[tbp]
\centerline{\includegraphics[width=0.95\columnwidth]{figs/fig5_discrimination.pdf}}
\caption{\textbf{Discrimination metric
$\Delta D(V) = d^2I_{\mathrm{mol}}/dV^2 - d^2I_{\mathrm{Baseline}}/dV^2$.}
Baseline-subtracted second-derivative spectra isolate the molecular IETS
contribution from the bulk-phonon and coherent backgrounds. The 100~meV
and 180~meV modes produce non-overlapping features separated by 80~mV,
well outside the thermal broadening $3k_BT/e \approx 78$~mV at 300~K. For
Mol\_AB the two features add linearly ($\Delta D_{AB}(V) \approx
\Delta D_A(V) + \Delta D_B(V)$ within numerical noise), a direct
demonstration of the closed-form analytic second-derivative (\ref{eq:d2i})
acting on a mixed analyte. The peak-to-peak values
$|\Delta D|_{\max}$ for each single-mode species are the figures of merit
for pixel-level discrimination. No lock-in or smoothing of any kind is
applied: the signal here is the direct output of the closed-form
(\ref{eq:d2i}) minus its Baseline evaluation.}
\label{fig:disc}
\end{figure}
.
%%
%% Each caption is written to be fully self-contained: an expert in the
%% field who reads only the paper title and this figure + caption should
%% be able to tell what was done, how, and what was shown.
%%
%% Figure files (to be generated from results/2026-04-09/*.npz by
%% run/plot_sispad_figs.py):
%%
%%   figs/fig1_device.pdf     — Device stack, band diagram, chi(z), rank-1 loop
%%   figs/fig2_patil.pdf      — Patil 1D benchmark + conservation residual inset
%%   figs/fig3_iv.pdf         — I–V of Baseline/Mol_A/Mol_B/Mol_AB on ZnO/MgZnO
%%   figs/fig4_d2i.pdf        — Analytic d2I/dV2 for all 4 molecules
%%   figs/fig5_discrimination.pdf — Baseline-subtracted discrimination metric
%%
%% Usage in sispad2026_abstract.tex:
%%   %% SISPAD 2026 — figures + captions, kept separate from the main abstract
%% file so captions can be iterated independently and dropped into
%% papers/sispad2026_abstract.tex via %% SISPAD 2026 — figures + captions, kept separate from the main abstract
%% file so captions can be iterated independently and dropped into
%% papers/sispad2026_abstract.tex via \input{sispad2026_figures.tex}.
%%
%% Each caption is written to be fully self-contained: an expert in the
%% field who reads only the paper title and this figure + caption should
%% be able to tell what was done, how, and what was shown.
%%
%% Figure files (to be generated from results/2026-04-09/*.npz by
%% run/plot_sispad_figs.py):
%%
%%   figs/fig1_device.pdf     — Device stack, band diagram, chi(z), rank-1 loop
%%   figs/fig2_patil.pdf      — Patil 1D benchmark + conservation residual inset
%%   figs/fig3_iv.pdf         — I–V of Baseline/Mol_A/Mol_B/Mol_AB on ZnO/MgZnO
%%   figs/fig4_d2i.pdf        — Analytic d2I/dV2 for all 4 molecules
%%   figs/fig5_discrimination.pdf — Baseline-subtracted discrimination metric
%%
%% Usage in sispad2026_abstract.tex:
%%   \input{sispad2026_figures}
%% or drop the individual \begin{figure}…\end{figure} blocks inline.
%% ----------------------------------------------------------------------

% --- Fig. 1 ----------------------------------------------------------------
\begin{figure}[tbp]
\centerline{\includegraphics[width=0.95\columnwidth]{figs/fig1_device.pdf}}
\caption{\textbf{Device and methodology at a glance.}
(a)~Layer stack of the asymmetric
ZnO/Mg\textsubscript{0.3}Zn\textsubscript{0.7}O RTD: 10~nm $n^+$-ZnO emitter
contact, 3~nm Mg\textsubscript{0.3}Zn\textsubscript{0.7}O emitter barrier
(host of the sensing molecule), 3~nm ZnO quantum well, 1.5~nm
Mg\textsubscript{0.3}Zn\textsubscript{0.7}O collector barrier, and 10~nm
$n^+$-ZnO collector. Conduction-band offset $\Delta E_c = 0.47$~eV,
effective mass $m^* = 0.28\,m_0$. The transverse area is
$A = 10\,\mu\mathrm{m} \times 10\,\mu\mathrm{m}$, entering the current
expression as a Datta ``Fix-$A$'' prefactor.
(b)~Conduction-band profile at $V = 0$ and at $V = V_{\mathrm{res}}$, showing
the quantized well level $E_{\mathrm{res}}$ crossing $E_F$ from above to
give negative differential resistance.
(c)~Molecular form factor $\chi(z)$, a Gaussian of width
$\sigma_{\mathrm{mol}} = 3$~\AA\ centered at the emitter-barrier midpoint,
defining the rank-1 electron--vibron coupling
$M = D_0\,\chi(z)\,|u\rangle\langle u|$.
(d)~Schematic of the rank-1 projected SCBA loop: a single
scalar-in-mode-space, 1D-in-$z$ Dyson equation dressed self-consistently by
the phonon self-energies $\tilde\sigma^{R,<,>}$. This replaces the full
nine-mode coupled SCBA with machine-precision accuracy whenever
$\sigma_{\mathrm{mol}} k_\perp \ll 1$.}
\label{fig:device}
\end{figure}

% --- Fig. 2 ----------------------------------------------------------------
\begin{figure}[tbp]
\centerline{\includegraphics[width=0.95\columnwidth]{figs/fig2_patil.pdf}}
\caption{\textbf{Validation against the Patil 1D benchmark (Tier-1 $\gamma$,
$\delta$).}
Right-contact current $I_R(V)$ of Patil's 1D symmetric RTD computed with
the rank-1 Keldysh wrapper (red circles) overlaid on the Patil 1D reference
from the published MATLAB code \texttt{rtd2modes\_1d.m} (black line). The
rank-1 solver reproduces the NDR peak to 0.55\%:
$I_R = 5.667 \times 10^{-10}$~A (rank-1) versus
$5.698 \times 10^{-10}$~A (reference) at $V = 0.336$~V. The orange
triangles show $-I_L$ on the same axes; the two traces would be exactly
coincident for a perfectly converged SCBA solution.
Inset: the per-bias current-conservation residual
$|I_L + I_R| / \max(|I_L|,|I_R|)$ in log scale. Exposing both contacts lets
us read $\delta$ directly from the Meir--Wingreen trace without any
auxiliary bookkeeping. All parameters (device, mesh, bias convention,
phonon modes, coupling, temperature) are identical to the reference.}
\label{fig:patil}
\end{figure}

% --- Fig. 3 ----------------------------------------------------------------
\begin{figure}[tbp]
\centerline{\includegraphics[width=0.95\columnwidth]{figs/fig3_iv.pdf}}
\caption{\textbf{Current--voltage characteristics of the ZnO/MgZnO primary
stack at 300~K.}
Rank-1 projected SCBA sweep (\texttt{max\_iter}$=10$, mix$=0.4$) on the
asymmetric ZnO/Mg\textsubscript{0.3}Zn\textsubscript{0.7}O RTD for the
four test species: Baseline (blank reference, bulk ZnO LO phonon at 72~meV
only, no molecule), Mol\_A (single vibrational mode $\hbar\omega_A =
100$~meV), Mol\_B ($\hbar\omega_B = 180$~meV), and Mol\_AB (both modes).
All molecules couple with $D_0^2 = 0.1\,\mathrm{eV}^2$ through a Gaussian
form factor of width 3~\AA\ centered on the emitter-barrier midpoint. Each
trace shows a clean NDR feature at $V \approx V_{\mathrm{res}}$; the phonon
sidebands pull current into the sub-resonance region, shifting the turn-on
and broadening the NDR step in a mode-specific way. Bias grid: 201~points
over 0--400~mV. Energy grid: $dE = 2$~meV over $[-0.25, +0.5]$~eV.
$E_F = 20$~meV. No numerical differentiation is applied to $I(V)$
anywhere in this work; the $d^2I/dV^2$ in Fig.~\ref{fig:d2i} comes from
the closed-form formula~(\ref{eq:d2i}) evaluated on the same converged
Keldysh state used to build this figure.}
\label{fig:iv}
\end{figure}

% --- Fig. 4 ----------------------------------------------------------------
\begin{figure}[tbp]
\centerline{\includegraphics[width=0.95\columnwidth]{figs/fig4_d2i.pdf}}
\caption{\textbf{Analytic $d^2I_R/dV^2$ --- vibrational fingerprints at
300~K.}
The second derivative is computed in closed form directly from the
converged Keldysh state at each bias via
$d^2I_R/dV^2 = (e^2/h)\cdot(1/4)\int dE\,[f_L''(E,V)\,T_{RL}(E) +
f_R''(E,V)\,(T_{RR}(E) - T_{RA}(E))]$, with the kernels
$T_{RL}, T_{RR}, T_{RA}$ built once from
$(G^R,\Gamma_L,\Gamma_R)$ at each bias --- no numerical $V$-grid
differentiation. At $D_0=0$ the formula reduces to the Landauer coherent
expression by the optical theorem
$T_{RA} = T_{RL} + T_{RR}$ (unit-tested to $10^{-6}$).
Mol\_A shows a peak at $V \approx V_{\mathrm{res}} + \hbar\omega_A/e \approx
187$~mV; Mol\_B at $V \approx V_{\mathrm{res}} + \hbar\omega_B/e \approx
267$~mV; Mol\_AB shows both features simultaneously. Baseline has a shared
feature at $\sim 162$~mV from the 72~meV bulk LO phonon. Red ticks mark the
predicted sideband biases: these are \emph{not fits}.}
\label{fig:d2i}
\end{figure}

% --- Fig. 5 ----------------------------------------------------------------
\begin{figure}[tbp]
\centerline{\includegraphics[width=0.95\columnwidth]{figs/fig5_discrimination.pdf}}
\caption{\textbf{Discrimination metric
$\Delta D(V) = d^2I_{\mathrm{mol}}/dV^2 - d^2I_{\mathrm{Baseline}}/dV^2$.}
Baseline-subtracted second-derivative spectra isolate the molecular IETS
contribution from the bulk-phonon and coherent backgrounds. The 100~meV
and 180~meV modes produce non-overlapping features separated by 80~mV,
well outside the thermal broadening $3k_BT/e \approx 78$~mV at 300~K. For
Mol\_AB the two features add linearly ($\Delta D_{AB}(V) \approx
\Delta D_A(V) + \Delta D_B(V)$ within numerical noise), a direct
demonstration of the closed-form analytic second-derivative (\ref{eq:d2i})
acting on a mixed analyte. The peak-to-peak values
$|\Delta D|_{\max}$ for each single-mode species are the figures of merit
for pixel-level discrimination. No lock-in or smoothing of any kind is
applied: the signal here is the direct output of the closed-form
(\ref{eq:d2i}) minus its Baseline evaluation.}
\label{fig:disc}
\end{figure}
.
%%
%% Each caption is written to be fully self-contained: an expert in the
%% field who reads only the paper title and this figure + caption should
%% be able to tell what was done, how, and what was shown.
%%
%% Figure files (to be generated from results/2026-04-09/*.npz by
%% run/plot_sispad_figs.py):
%%
%%   figs/fig1_device.pdf     — Device stack, band diagram, chi(z), rank-1 loop
%%   figs/fig2_patil.pdf      — Patil 1D benchmark + conservation residual inset
%%   figs/fig3_iv.pdf         — I–V of Baseline/Mol_A/Mol_B/Mol_AB on ZnO/MgZnO
%%   figs/fig4_d2i.pdf        — Analytic d2I/dV2 for all 4 molecules
%%   figs/fig5_discrimination.pdf — Baseline-subtracted discrimination metric
%%
%% Usage in sispad2026_abstract.tex:
%%   %% SISPAD 2026 — figures + captions, kept separate from the main abstract
%% file so captions can be iterated independently and dropped into
%% papers/sispad2026_abstract.tex via \input{sispad2026_figures.tex}.
%%
%% Each caption is written to be fully self-contained: an expert in the
%% field who reads only the paper title and this figure + caption should
%% be able to tell what was done, how, and what was shown.
%%
%% Figure files (to be generated from results/2026-04-09/*.npz by
%% run/plot_sispad_figs.py):
%%
%%   figs/fig1_device.pdf     — Device stack, band diagram, chi(z), rank-1 loop
%%   figs/fig2_patil.pdf      — Patil 1D benchmark + conservation residual inset
%%   figs/fig3_iv.pdf         — I–V of Baseline/Mol_A/Mol_B/Mol_AB on ZnO/MgZnO
%%   figs/fig4_d2i.pdf        — Analytic d2I/dV2 for all 4 molecules
%%   figs/fig5_discrimination.pdf — Baseline-subtracted discrimination metric
%%
%% Usage in sispad2026_abstract.tex:
%%   \input{sispad2026_figures}
%% or drop the individual \begin{figure}…\end{figure} blocks inline.
%% ----------------------------------------------------------------------

% --- Fig. 1 ----------------------------------------------------------------
\begin{figure}[tbp]
\centerline{\includegraphics[width=0.95\columnwidth]{figs/fig1_device.pdf}}
\caption{\textbf{Device and methodology at a glance.}
(a)~Layer stack of the asymmetric
ZnO/Mg\textsubscript{0.3}Zn\textsubscript{0.7}O RTD: 10~nm $n^+$-ZnO emitter
contact, 3~nm Mg\textsubscript{0.3}Zn\textsubscript{0.7}O emitter barrier
(host of the sensing molecule), 3~nm ZnO quantum well, 1.5~nm
Mg\textsubscript{0.3}Zn\textsubscript{0.7}O collector barrier, and 10~nm
$n^+$-ZnO collector. Conduction-band offset $\Delta E_c = 0.47$~eV,
effective mass $m^* = 0.28\,m_0$. The transverse area is
$A = 10\,\mu\mathrm{m} \times 10\,\mu\mathrm{m}$, entering the current
expression as a Datta ``Fix-$A$'' prefactor.
(b)~Conduction-band profile at $V = 0$ and at $V = V_{\mathrm{res}}$, showing
the quantized well level $E_{\mathrm{res}}$ crossing $E_F$ from above to
give negative differential resistance.
(c)~Molecular form factor $\chi(z)$, a Gaussian of width
$\sigma_{\mathrm{mol}} = 3$~\AA\ centered at the emitter-barrier midpoint,
defining the rank-1 electron--vibron coupling
$M = D_0\,\chi(z)\,|u\rangle\langle u|$.
(d)~Schematic of the rank-1 projected SCBA loop: a single
scalar-in-mode-space, 1D-in-$z$ Dyson equation dressed self-consistently by
the phonon self-energies $\tilde\sigma^{R,<,>}$. This replaces the full
nine-mode coupled SCBA with machine-precision accuracy whenever
$\sigma_{\mathrm{mol}} k_\perp \ll 1$.}
\label{fig:device}
\end{figure}

% --- Fig. 2 ----------------------------------------------------------------
\begin{figure}[tbp]
\centerline{\includegraphics[width=0.95\columnwidth]{figs/fig2_patil.pdf}}
\caption{\textbf{Validation against the Patil 1D benchmark (Tier-1 $\gamma$,
$\delta$).}
Right-contact current $I_R(V)$ of Patil's 1D symmetric RTD computed with
the rank-1 Keldysh wrapper (red circles) overlaid on the Patil 1D reference
from the published MATLAB code \texttt{rtd2modes\_1d.m} (black line). The
rank-1 solver reproduces the NDR peak to 0.55\%:
$I_R = 5.667 \times 10^{-10}$~A (rank-1) versus
$5.698 \times 10^{-10}$~A (reference) at $V = 0.336$~V. The orange
triangles show $-I_L$ on the same axes; the two traces would be exactly
coincident for a perfectly converged SCBA solution.
Inset: the per-bias current-conservation residual
$|I_L + I_R| / \max(|I_L|,|I_R|)$ in log scale. Exposing both contacts lets
us read $\delta$ directly from the Meir--Wingreen trace without any
auxiliary bookkeeping. All parameters (device, mesh, bias convention,
phonon modes, coupling, temperature) are identical to the reference.}
\label{fig:patil}
\end{figure}

% --- Fig. 3 ----------------------------------------------------------------
\begin{figure}[tbp]
\centerline{\includegraphics[width=0.95\columnwidth]{figs/fig3_iv.pdf}}
\caption{\textbf{Current--voltage characteristics of the ZnO/MgZnO primary
stack at 300~K.}
Rank-1 projected SCBA sweep (\texttt{max\_iter}$=10$, mix$=0.4$) on the
asymmetric ZnO/Mg\textsubscript{0.3}Zn\textsubscript{0.7}O RTD for the
four test species: Baseline (blank reference, bulk ZnO LO phonon at 72~meV
only, no molecule), Mol\_A (single vibrational mode $\hbar\omega_A =
100$~meV), Mol\_B ($\hbar\omega_B = 180$~meV), and Mol\_AB (both modes).
All molecules couple with $D_0^2 = 0.1\,\mathrm{eV}^2$ through a Gaussian
form factor of width 3~\AA\ centered on the emitter-barrier midpoint. Each
trace shows a clean NDR feature at $V \approx V_{\mathrm{res}}$; the phonon
sidebands pull current into the sub-resonance region, shifting the turn-on
and broadening the NDR step in a mode-specific way. Bias grid: 201~points
over 0--400~mV. Energy grid: $dE = 2$~meV over $[-0.25, +0.5]$~eV.
$E_F = 20$~meV. No numerical differentiation is applied to $I(V)$
anywhere in this work; the $d^2I/dV^2$ in Fig.~\ref{fig:d2i} comes from
the closed-form formula~(\ref{eq:d2i}) evaluated on the same converged
Keldysh state used to build this figure.}
\label{fig:iv}
\end{figure}

% --- Fig. 4 ----------------------------------------------------------------
\begin{figure}[tbp]
\centerline{\includegraphics[width=0.95\columnwidth]{figs/fig4_d2i.pdf}}
\caption{\textbf{Analytic $d^2I_R/dV^2$ --- vibrational fingerprints at
300~K.}
The second derivative is computed in closed form directly from the
converged Keldysh state at each bias via
$d^2I_R/dV^2 = (e^2/h)\cdot(1/4)\int dE\,[f_L''(E,V)\,T_{RL}(E) +
f_R''(E,V)\,(T_{RR}(E) - T_{RA}(E))]$, with the kernels
$T_{RL}, T_{RR}, T_{RA}$ built once from
$(G^R,\Gamma_L,\Gamma_R)$ at each bias --- no numerical $V$-grid
differentiation. At $D_0=0$ the formula reduces to the Landauer coherent
expression by the optical theorem
$T_{RA} = T_{RL} + T_{RR}$ (unit-tested to $10^{-6}$).
Mol\_A shows a peak at $V \approx V_{\mathrm{res}} + \hbar\omega_A/e \approx
187$~mV; Mol\_B at $V \approx V_{\mathrm{res}} + \hbar\omega_B/e \approx
267$~mV; Mol\_AB shows both features simultaneously. Baseline has a shared
feature at $\sim 162$~mV from the 72~meV bulk LO phonon. Red ticks mark the
predicted sideband biases: these are \emph{not fits}.}
\label{fig:d2i}
\end{figure}

% --- Fig. 5 ----------------------------------------------------------------
\begin{figure}[tbp]
\centerline{\includegraphics[width=0.95\columnwidth]{figs/fig5_discrimination.pdf}}
\caption{\textbf{Discrimination metric
$\Delta D(V) = d^2I_{\mathrm{mol}}/dV^2 - d^2I_{\mathrm{Baseline}}/dV^2$.}
Baseline-subtracted second-derivative spectra isolate the molecular IETS
contribution from the bulk-phonon and coherent backgrounds. The 100~meV
and 180~meV modes produce non-overlapping features separated by 80~mV,
well outside the thermal broadening $3k_BT/e \approx 78$~mV at 300~K. For
Mol\_AB the two features add linearly ($\Delta D_{AB}(V) \approx
\Delta D_A(V) + \Delta D_B(V)$ within numerical noise), a direct
demonstration of the closed-form analytic second-derivative (\ref{eq:d2i})
acting on a mixed analyte. The peak-to-peak values
$|\Delta D|_{\max}$ for each single-mode species are the figures of merit
for pixel-level discrimination. No lock-in or smoothing of any kind is
applied: the signal here is the direct output of the closed-form
(\ref{eq:d2i}) minus its Baseline evaluation.}
\label{fig:disc}
\end{figure}

%% or drop the individual \begin{figure}…\end{figure} blocks inline.
%% ----------------------------------------------------------------------

% --- Fig. 1 ----------------------------------------------------------------
\begin{figure}[tbp]
\centerline{\includegraphics[width=0.95\columnwidth]{figs/fig1_device.pdf}}
\caption{\textbf{Device and methodology at a glance.}
(a)~Layer stack of the asymmetric
ZnO/Mg\textsubscript{0.3}Zn\textsubscript{0.7}O RTD: 10~nm $n^+$-ZnO emitter
contact, 3~nm Mg\textsubscript{0.3}Zn\textsubscript{0.7}O emitter barrier
(host of the sensing molecule), 3~nm ZnO quantum well, 1.5~nm
Mg\textsubscript{0.3}Zn\textsubscript{0.7}O collector barrier, and 10~nm
$n^+$-ZnO collector. Conduction-band offset $\Delta E_c = 0.47$~eV,
effective mass $m^* = 0.28\,m_0$. The transverse area is
$A = 10\,\mu\mathrm{m} \times 10\,\mu\mathrm{m}$, entering the current
expression as a Datta ``Fix-$A$'' prefactor.
(b)~Conduction-band profile at $V = 0$ and at $V = V_{\mathrm{res}}$, showing
the quantized well level $E_{\mathrm{res}}$ crossing $E_F$ from above to
give negative differential resistance.
(c)~Molecular form factor $\chi(z)$, a Gaussian of width
$\sigma_{\mathrm{mol}} = 3$~\AA\ centered at the emitter-barrier midpoint,
defining the rank-1 electron--vibron coupling
$M = D_0\,\chi(z)\,|u\rangle\langle u|$.
(d)~Schematic of the rank-1 projected SCBA loop: a single
scalar-in-mode-space, 1D-in-$z$ Dyson equation dressed self-consistently by
the phonon self-energies $\tilde\sigma^{R,<,>}$. This replaces the full
nine-mode coupled SCBA with machine-precision accuracy whenever
$\sigma_{\mathrm{mol}} k_\perp \ll 1$.}
\label{fig:device}
\end{figure}

% --- Fig. 2 ----------------------------------------------------------------
\begin{figure}[tbp]
\centerline{\includegraphics[width=0.95\columnwidth]{figs/fig2_patil.pdf}}
\caption{\textbf{Validation against the Patil 1D benchmark (Tier-1 $\gamma$,
$\delta$).}
Right-contact current $I_R(V)$ of Patil's 1D symmetric RTD computed with
the rank-1 Keldysh wrapper (red circles) overlaid on the Patil 1D reference
from the published MATLAB code \texttt{rtd2modes\_1d.m} (black line). The
rank-1 solver reproduces the NDR peak to 0.55\%:
$I_R = 5.667 \times 10^{-10}$~A (rank-1) versus
$5.698 \times 10^{-10}$~A (reference) at $V = 0.336$~V. The orange
triangles show $-I_L$ on the same axes; the two traces would be exactly
coincident for a perfectly converged SCBA solution.
Inset: the per-bias current-conservation residual
$|I_L + I_R| / \max(|I_L|,|I_R|)$ in log scale. Exposing both contacts lets
us read $\delta$ directly from the Meir--Wingreen trace without any
auxiliary bookkeeping. All parameters (device, mesh, bias convention,
phonon modes, coupling, temperature) are identical to the reference.}
\label{fig:patil}
\end{figure}

% --- Fig. 3 ----------------------------------------------------------------
\begin{figure}[tbp]
\centerline{\includegraphics[width=0.95\columnwidth]{figs/fig3_iv.pdf}}
\caption{\textbf{Current--voltage characteristics of the ZnO/MgZnO primary
stack at 300~K.}
Rank-1 projected SCBA sweep (\texttt{max\_iter}$=10$, mix$=0.4$) on the
asymmetric ZnO/Mg\textsubscript{0.3}Zn\textsubscript{0.7}O RTD for the
four test species: Baseline (blank reference, bulk ZnO LO phonon at 72~meV
only, no molecule), Mol\_A (single vibrational mode $\hbar\omega_A =
100$~meV), Mol\_B ($\hbar\omega_B = 180$~meV), and Mol\_AB (both modes).
All molecules couple with $D_0^2 = 0.1\,\mathrm{eV}^2$ through a Gaussian
form factor of width 3~\AA\ centered on the emitter-barrier midpoint. Each
trace shows a clean NDR feature at $V \approx V_{\mathrm{res}}$; the phonon
sidebands pull current into the sub-resonance region, shifting the turn-on
and broadening the NDR step in a mode-specific way. Bias grid: 201~points
over 0--400~mV. Energy grid: $dE = 2$~meV over $[-0.25, +0.5]$~eV.
$E_F = 20$~meV. No numerical differentiation is applied to $I(V)$
anywhere in this work; the $d^2I/dV^2$ in Fig.~\ref{fig:d2i} comes from
the closed-form formula~(\ref{eq:d2i}) evaluated on the same converged
Keldysh state used to build this figure.}
\label{fig:iv}
\end{figure}

% --- Fig. 4 ----------------------------------------------------------------
\begin{figure}[tbp]
\centerline{\includegraphics[width=0.95\columnwidth]{figs/fig4_d2i.pdf}}
\caption{\textbf{Analytic $d^2I_R/dV^2$ --- vibrational fingerprints at
300~K.}
The second derivative is computed in closed form directly from the
converged Keldysh state at each bias via
$d^2I_R/dV^2 = (e^2/h)\cdot(1/4)\int dE\,[f_L''(E,V)\,T_{RL}(E) +
f_R''(E,V)\,(T_{RR}(E) - T_{RA}(E))]$, with the kernels
$T_{RL}, T_{RR}, T_{RA}$ built once from
$(G^R,\Gamma_L,\Gamma_R)$ at each bias --- no numerical $V$-grid
differentiation. At $D_0=0$ the formula reduces to the Landauer coherent
expression by the optical theorem
$T_{RA} = T_{RL} + T_{RR}$ (unit-tested to $10^{-6}$).
Mol\_A shows a peak at $V \approx V_{\mathrm{res}} + \hbar\omega_A/e \approx
187$~mV; Mol\_B at $V \approx V_{\mathrm{res}} + \hbar\omega_B/e \approx
267$~mV; Mol\_AB shows both features simultaneously. Baseline has a shared
feature at $\sim 162$~mV from the 72~meV bulk LO phonon. Red ticks mark the
predicted sideband biases: these are \emph{not fits}.}
\label{fig:d2i}
\end{figure}

% --- Fig. 5 ----------------------------------------------------------------
\begin{figure}[tbp]
\centerline{\includegraphics[width=0.95\columnwidth]{figs/fig5_discrimination.pdf}}
\caption{\textbf{Discrimination metric
$\Delta D(V) = d^2I_{\mathrm{mol}}/dV^2 - d^2I_{\mathrm{Baseline}}/dV^2$.}
Baseline-subtracted second-derivative spectra isolate the molecular IETS
contribution from the bulk-phonon and coherent backgrounds. The 100~meV
and 180~meV modes produce non-overlapping features separated by 80~mV,
well outside the thermal broadening $3k_BT/e \approx 78$~mV at 300~K. For
Mol\_AB the two features add linearly ($\Delta D_{AB}(V) \approx
\Delta D_A(V) + \Delta D_B(V)$ within numerical noise), a direct
demonstration of the closed-form analytic second-derivative (\ref{eq:d2i})
acting on a mixed analyte. The peak-to-peak values
$|\Delta D|_{\max}$ for each single-mode species are the figures of merit
for pixel-level discrimination. No lock-in or smoothing of any kind is
applied: the signal here is the direct output of the closed-form
(\ref{eq:d2i}) minus its Baseline evaluation.}
\label{fig:disc}
\end{figure}

%% or drop the individual \begin{figure}…\end{figure} blocks inline.
%% ----------------------------------------------------------------------

% --- Fig. 1 ----------------------------------------------------------------
\begin{figure}[tbp]
\centerline{\includegraphics[width=0.95\columnwidth]{figs/fig1_device.pdf}}
\caption{\textbf{Device and methodology at a glance.}
(a)~Layer stack of the asymmetric
ZnO/Mg\textsubscript{0.3}Zn\textsubscript{0.7}O RTD: 10~nm $n^+$-ZnO emitter
contact, 3~nm Mg\textsubscript{0.3}Zn\textsubscript{0.7}O emitter barrier
(host of the sensing molecule), 3~nm ZnO quantum well, 1.5~nm
Mg\textsubscript{0.3}Zn\textsubscript{0.7}O collector barrier, and 10~nm
$n^+$-ZnO collector. Conduction-band offset $\Delta E_c = 0.47$~eV,
effective mass $m^* = 0.28\,m_0$. The transverse area is
$A = 10\,\mu\mathrm{m} \times 10\,\mu\mathrm{m}$, entering the current
expression as a Datta ``Fix-$A$'' prefactor.
(b)~Conduction-band profile at $V = 0$ and at $V = V_{\mathrm{res}}$, showing
the quantized well level $E_{\mathrm{res}}$ crossing $E_F$ from above to
give negative differential resistance.
(c)~Molecular form factor $\chi(z)$, a Gaussian of width
$\sigma_{\mathrm{mol}} = 3$~\AA\ centered at the emitter-barrier midpoint,
defining the rank-1 electron--vibron coupling
$M = D_0\,\chi(z)\,|u\rangle\langle u|$.
(d)~Schematic of the rank-1 projected SCBA loop: a single
scalar-in-mode-space, 1D-in-$z$ Dyson equation dressed self-consistently by
the phonon self-energies $\tilde\sigma^{R,<,>}$. This replaces the full
nine-mode coupled SCBA with machine-precision accuracy whenever
$\sigma_{\mathrm{mol}} k_\perp \ll 1$.}
\label{fig:device}
\end{figure}

% --- Fig. 2 ----------------------------------------------------------------
\begin{figure}[tbp]
\centerline{\includegraphics[width=0.95\columnwidth]{figs/fig2_patil.pdf}}
\caption{\textbf{Validation against the Patil 1D benchmark (Tier-1 $\gamma$,
$\delta$).}
Right-contact current $I_R(V)$ of Patil's 1D symmetric RTD computed with
the rank-1 Keldysh wrapper (red circles) overlaid on the Patil 1D reference
from the published MATLAB code \texttt{rtd2modes\_1d.m} (black line). The
rank-1 solver reproduces the NDR peak to 0.55\%:
$I_R = 5.667 \times 10^{-10}$~A (rank-1) versus
$5.698 \times 10^{-10}$~A (reference) at $V = 0.336$~V. The orange
triangles show $-I_L$ on the same axes; the two traces would be exactly
coincident for a perfectly converged SCBA solution.
Inset: the per-bias current-conservation residual
$|I_L + I_R| / \max(|I_L|,|I_R|)$ in log scale. Exposing both contacts lets
us read $\delta$ directly from the Meir--Wingreen trace without any
auxiliary bookkeeping. All parameters (device, mesh, bias convention,
phonon modes, coupling, temperature) are identical to the reference.}
\label{fig:patil}
\end{figure}

% --- Fig. 3 ----------------------------------------------------------------
\begin{figure}[tbp]
\centerline{\includegraphics[width=0.95\columnwidth]{figs/fig3_iv.pdf}}
\caption{\textbf{Current--voltage characteristics of the ZnO/MgZnO primary
stack at 300~K.}
Rank-1 projected SCBA sweep (\texttt{max\_iter}$=10$, mix$=0.4$) on the
asymmetric ZnO/Mg\textsubscript{0.3}Zn\textsubscript{0.7}O RTD for the
four test species: Baseline (blank reference, bulk ZnO LO phonon at 72~meV
only, no molecule), Mol\_A (single vibrational mode $\hbar\omega_A =
100$~meV), Mol\_B ($\hbar\omega_B = 180$~meV), and Mol\_AB (both modes).
All molecules couple with $D_0^2 = 0.1\,\mathrm{eV}^2$ through a Gaussian
form factor of width 3~\AA\ centered on the emitter-barrier midpoint. Each
trace shows a clean NDR feature at $V \approx V_{\mathrm{res}}$; the phonon
sidebands pull current into the sub-resonance region, shifting the turn-on
and broadening the NDR step in a mode-specific way. Bias grid: 201~points
over 0--400~mV. Energy grid: $dE = 2$~meV over $[-0.25, +0.5]$~eV.
$E_F = 20$~meV. No numerical differentiation is applied to $I(V)$
anywhere in this work; the $d^2I/dV^2$ in Fig.~\ref{fig:d2i} comes from
the closed-form formula~(\ref{eq:d2i}) evaluated on the same converged
Keldysh state used to build this figure.}
\label{fig:iv}
\end{figure}

% --- Fig. 4 ----------------------------------------------------------------
\begin{figure}[tbp]
\centerline{\includegraphics[width=0.95\columnwidth]{figs/fig4_d2i.pdf}}
\caption{\textbf{Analytic $d^2I_R/dV^2$ --- vibrational fingerprints at
300~K.}
The second derivative is computed in closed form directly from the
converged Keldysh state at each bias via
$d^2I_R/dV^2 = (e^2/h)\cdot(1/4)\int dE\,[f_L''(E,V)\,T_{RL}(E) +
f_R''(E,V)\,(T_{RR}(E) - T_{RA}(E))]$, with the kernels
$T_{RL}, T_{RR}, T_{RA}$ built once from
$(G^R,\Gamma_L,\Gamma_R)$ at each bias --- no numerical $V$-grid
differentiation. At $D_0=0$ the formula reduces to the Landauer coherent
expression by the optical theorem
$T_{RA} = T_{RL} + T_{RR}$ (unit-tested to $10^{-6}$).
Mol\_A shows a peak at $V \approx V_{\mathrm{res}} + \hbar\omega_A/e \approx
187$~mV; Mol\_B at $V \approx V_{\mathrm{res}} + \hbar\omega_B/e \approx
267$~mV; Mol\_AB shows both features simultaneously. Baseline has a shared
feature at $\sim 162$~mV from the 72~meV bulk LO phonon. Red ticks mark the
predicted sideband biases: these are \emph{not fits}.}
\label{fig:d2i}
\end{figure}

% --- Fig. 5 ----------------------------------------------------------------
\begin{figure}[tbp]
\centerline{\includegraphics[width=0.95\columnwidth]{figs/fig5_discrimination.pdf}}
\caption{\textbf{Discrimination metric
$\Delta D(V) = d^2I_{\mathrm{mol}}/dV^2 - d^2I_{\mathrm{Baseline}}/dV^2$.}
Baseline-subtracted second-derivative spectra isolate the molecular IETS
contribution from the bulk-phonon and coherent backgrounds. The 100~meV
and 180~meV modes produce non-overlapping features separated by 80~mV,
well outside the thermal broadening $3k_BT/e \approx 78$~mV at 300~K. For
Mol\_AB the two features add linearly ($\Delta D_{AB}(V) \approx
\Delta D_A(V) + \Delta D_B(V)$ within numerical noise), a direct
demonstration of the closed-form analytic second-derivative (\ref{eq:d2i})
acting on a mixed analyte. The peak-to-peak values
$|\Delta D|_{\max}$ for each single-mode species are the figures of merit
for pixel-level discrimination. No lock-in or smoothing of any kind is
applied: the signal here is the direct output of the closed-form
(\ref{eq:d2i}) minus its Baseline evaluation.}
\label{fig:disc}
\end{figure}

%% or drop the individual \begin{figure}…\end{figure} blocks inline.
%% ----------------------------------------------------------------------

% --- Fig. 1 ----------------------------------------------------------------
\begin{figure}[tbp]
\centerline{\includegraphics[width=0.95\columnwidth]{figs/fig1_device.pdf}}
\caption{\textbf{Device and methodology at a glance.}
(a)~Layer stack of the asymmetric
ZnO/Mg\textsubscript{0.3}Zn\textsubscript{0.7}O RTD: 10~nm $n^+$-ZnO emitter
contact, 3~nm Mg\textsubscript{0.3}Zn\textsubscript{0.7}O emitter barrier
(host of the sensing molecule), 3~nm ZnO quantum well, 1.5~nm
Mg\textsubscript{0.3}Zn\textsubscript{0.7}O collector barrier, and 10~nm
$n^+$-ZnO collector. Conduction-band offset $\Delta E_c = 0.47$~eV,
effective mass $m^* = 0.28\,m_0$. The transverse area is
$A = 10\,\mu\mathrm{m} \times 10\,\mu\mathrm{m}$, entering the current
expression as a Datta ``Fix-$A$'' prefactor.
(b)~Conduction-band profile at $V = 0$ and at $V = V_{\mathrm{res}}$, showing
the quantized well level $E_{\mathrm{res}}$ crossing $E_F$ from above to
give negative differential resistance.
(c)~Molecular form factor $\chi(z)$, a Gaussian of width
$\sigma_{\mathrm{mol}} = 3$~\AA\ centered at the emitter-barrier midpoint,
defining the rank-1 electron--vibron coupling
$M = D_0\,\chi(z)\,|u\rangle\langle u|$.
(d)~Schematic of the rank-1 projected SCBA loop: a single
scalar-in-mode-space, 1D-in-$z$ Dyson equation dressed self-consistently by
the phonon self-energies $\tilde\sigma^{R,<,>}$. This replaces the full
nine-mode coupled SCBA with machine-precision accuracy whenever
$\sigma_{\mathrm{mol}} k_\perp \ll 1$.}
\label{fig:device}
\end{figure}

% --- Fig. 2 ----------------------------------------------------------------
\begin{figure}[tbp]
\centerline{\includegraphics[width=0.95\columnwidth]{figs/fig2_patil.pdf}}
\caption{\textbf{Validation against the Patil 1D benchmark (Tier-1 $\gamma$,
$\delta$).}
Right-contact current $I_R(V)$ of Patil's 1D symmetric RTD computed with
the rank-1 Keldysh wrapper (red circles) overlaid on the Patil 1D reference
from the published MATLAB code \texttt{rtd2modes\_1d.m} (black line). The
rank-1 solver reproduces the NDR peak to 0.55\%:
$I_R = 5.667 \times 10^{-10}$~A (rank-1) versus
$5.698 \times 10^{-10}$~A (reference) at $V = 0.336$~V. The orange
triangles show $-I_L$ on the same axes; the two traces would be exactly
coincident for a perfectly converged SCBA solution.
Inset: the per-bias current-conservation residual
$|I_L + I_R| / \max(|I_L|,|I_R|)$ in log scale. Exposing both contacts lets
us read $\delta$ directly from the Meir--Wingreen trace without any
auxiliary bookkeeping. All parameters (device, mesh, bias convention,
phonon modes, coupling, temperature) are identical to the reference.}
\label{fig:patil}
\end{figure}

% --- Fig. 3 ----------------------------------------------------------------
\begin{figure}[tbp]
\centerline{\includegraphics[width=0.95\columnwidth]{figs/fig3_iv.pdf}}
\caption{\textbf{Current--voltage characteristics of the ZnO/MgZnO primary
stack at 300~K.}
Rank-1 projected SCBA sweep (\texttt{max\_iter}$=10$, mix$=0.4$) on the
asymmetric ZnO/Mg\textsubscript{0.3}Zn\textsubscript{0.7}O RTD for the
four test species: Baseline (blank reference, bulk ZnO LO phonon at 72~meV
only, no molecule), Mol\_A (single vibrational mode $\hbar\omega_A =
100$~meV), Mol\_B ($\hbar\omega_B = 180$~meV), and Mol\_AB (both modes).
All molecules couple with $D_0^2 = 0.1\,\mathrm{eV}^2$ through a Gaussian
form factor of width 3~\AA\ centered on the emitter-barrier midpoint. Each
trace shows a clean NDR feature at $V \approx V_{\mathrm{res}}$; the phonon
sidebands pull current into the sub-resonance region, shifting the turn-on
and broadening the NDR step in a mode-specific way. Bias grid: 201~points
over 0--400~mV. Energy grid: $dE = 2$~meV over $[-0.25, +0.5]$~eV.
$E_F = 20$~meV. No numerical differentiation is applied to $I(V)$
anywhere in this work; the $d^2I/dV^2$ in Fig.~\ref{fig:d2i} comes from
the closed-form formula~(\ref{eq:d2i}) evaluated on the same converged
Keldysh state used to build this figure.}
\label{fig:iv}
\end{figure}

% --- Fig. 4 ----------------------------------------------------------------
\begin{figure}[tbp]
\centerline{\includegraphics[width=0.95\columnwidth]{figs/fig4_d2i.pdf}}
\caption{\textbf{Analytic $d^2I_R/dV^2$ --- vibrational fingerprints at
300~K.}
The second derivative is computed in closed form directly from the
converged Keldysh state at each bias via
$d^2I_R/dV^2 = (e^2/h)\cdot(1/4)\int dE\,[f_L''(E,V)\,T_{RL}(E) +
f_R''(E,V)\,(T_{RR}(E) - T_{RA}(E))]$, with the kernels
$T_{RL}, T_{RR}, T_{RA}$ built once from
$(G^R,\Gamma_L,\Gamma_R)$ at each bias --- no numerical $V$-grid
differentiation. At $D_0=0$ the formula reduces to the Landauer coherent
expression by the optical theorem
$T_{RA} = T_{RL} + T_{RR}$ (unit-tested to $10^{-6}$).
Mol\_A shows a peak at $V \approx V_{\mathrm{res}} + \hbar\omega_A/e \approx
187$~mV; Mol\_B at $V \approx V_{\mathrm{res}} + \hbar\omega_B/e \approx
267$~mV; Mol\_AB shows both features simultaneously. Baseline has a shared
feature at $\sim 162$~mV from the 72~meV bulk LO phonon. Red ticks mark the
predicted sideband biases: these are \emph{not fits}.}
\label{fig:d2i}
\end{figure}

% --- Fig. 5 ----------------------------------------------------------------
\begin{figure}[tbp]
\centerline{\includegraphics[width=0.95\columnwidth]{figs/fig5_discrimination.pdf}}
\caption{\textbf{Discrimination metric
$\Delta D(V) = d^2I_{\mathrm{mol}}/dV^2 - d^2I_{\mathrm{Baseline}}/dV^2$.}
Baseline-subtracted second-derivative spectra isolate the molecular IETS
contribution from the bulk-phonon and coherent backgrounds. The 100~meV
and 180~meV modes produce non-overlapping features separated by 80~mV,
well outside the thermal broadening $3k_BT/e \approx 78$~mV at 300~K. For
Mol\_AB the two features add linearly ($\Delta D_{AB}(V) \approx
\Delta D_A(V) + \Delta D_B(V)$ within numerical noise), a direct
demonstration of the closed-form analytic second-derivative (\ref{eq:d2i})
acting on a mixed analyte. The peak-to-peak values
$|\Delta D|_{\max}$ for each single-mode species are the figures of merit
for pixel-level discrimination. No lock-in or smoothing of any kind is
applied: the signal here is the direct output of the closed-form
(\ref{eq:d2i}) minus its Baseline evaluation.}
\label{fig:disc}
\end{figure}
